\chapter{Proof of the non-collapsing}
\label{chap:proof-noncollapsing}
\label{sectnoncoll}

The precise statement of the non-collapsing result is given in the
next section. Essentially, the proof of non-collapsing in the
context of Ricci flow with surgery is the same as the proof in the
case of ordinary Ricci flows. Given a point $x\in {\mathcal M}$, one
finds a parabolic neighborhood whose size, $r'$, is determined by
the constants $r_i$, $C$ and $\epsilon$, contained in ${\bf
t}^{-1}([T_{i-1},T_i))$ and on which the curvature is bounded by
$(r')^{-2}$. Hence, by the inductive hypothesis, the final
time-slice of this neighborhood is $\kappa_i$-non-collapsed.
Furthermore, we can choose this neighborhood so that the reduced
${\mathcal L}$-length of its central point from $x$ is bounded by
$3/2$. This allows us to produce an open subset at an earlier time
whose reduced volume is bounded away from zero. Then using
Theorem~\ref{THM} we transfer this conclusion to a non-collapsing
conclusion for $x$. The main issue in this argument is to show that
there is a point in each earlier time-slice whose reduced length
from $x$ is at most $3/2$. We can argue as in the case of a Ricci
flow if we can show that any curve parameterized by backward time
starting at $x$ (a point where the hypothesis of
$\kappa$-non-collapsing holds) that comes close to  a surgery cap
either from above or below must have large ${\mathcal L}$-length. In
establishing the relevant estimates we are forced to require that
$\delta_i$ be sufficiently small.


\section{The statement of the non-collapsing result}

Here, we shall assume that after surgery the strong canonical
neighborhood assumption holds, and we shall establish the
non-collapsing result.

\begin{proposition}[Non-collapsing result assuming the strong canonical neighborhood assumption]
\label{prop:kappa-limit-noncollapsing}
\uses{thm:main-existence-ricci-flow-surgery}\label{KAPPALIMIT}
Suppose that for some $i\ge 0$ we have surgery parameter sequences $\delta_0\ge
\delta_1\ge\cdots\ge \delta_i>0$, $\epsilon=r_0\ge r_1\ge \cdots\ge
r_i>0$ and $\kappa_0\ge \kappa_1\ge \cdots\ge \kappa_i>0$.
 Then there is $0<\kappa\le \kappa_i$ and for any $0<r_{i+1}\le r_i$ there is
$0<\delta(r_{i+1})\le \delta_i$ such that the following holds. Suppose that
$\overline\delta\colon [0,T_{i+1}]\to \Ar^+$ is a non-increasing function with
$\overline\delta(t)\le \delta_j$ for all $t\in [T_j,T_{j+1})$ and
$\overline\delta(t)\le \delta(r_{i+1})$ for all $t\in [T_{i-1},T_{i+1})$.
Suppose that $({\mathcal M},G)$ is a Ricci flow with surgery defined for $0\le
t<T$ for some $T\in (T_i,T_{i+1}]$ with surgery control parameter
$\overline\delta(t)$. Suppose that the restriction of this Ricci flow with
surgery  to the time-interval $[0,T_i)$ satisfies the hypothesis of
Theorem~\ref{MAIN} with respect to the given sequences. Suppose also that the
entire Ricci flow with surgery $({\mathcal M},G)$ satisfies Assumptions (1) --
(7) and the strong $(C,\epsilon)$-canonical neighborhood assumption with
parameter $r_{i+1}$. Then $({\mathcal M},G)$ is $\kappa$-non-collapsed on all
scales $\le \epsilon$.
\end{proposition}

\begin{remark}[Dependence of the non-collapsing constants and the two-outcome alternative]
\label{rem:kappa-delta-dependence}
\uses{prop:kappa-limit-noncollapsing}
Implicitly,
$\kappa$ and $\delta(r_{i+1})$ are also allowed to depend on $t_0, \epsilon$,
and $C$, which are fixed, and also $i+1$.
Also recall that the non-collapsing condition allows for two outcomes: if $x$ is a point at which the
hypothesis of the non-collapsing hold, then
there is a lower bound on the volume of a ball centered at $x$, or  $x$
is contained in a component of its time-slice that has positive sectional curvature.
\end{remark}


\section{The proof of non-collapsing when $R(x)=r^{-2}$ with
$r\le r_{i+1}$}\label{rsmall}

Let us begin with an easy case of the non-collapsing result, where
non-collapsing follows easily from the strong canonical neighborhood
assumption, rather than from using ${\mathcal L}$-length and
monotonicity along ${\mathcal L}$-geodesics. We suppose that we have
a Ricci flow with surgery $({\mathcal M},G)$ defined for $0\le t<T$
with $T\in [T_i,T_{i+1})$, and a constant $r_{i+1}\le r_i$, all
satisfying the hypothesis of Proposition~\ref{KAPPALIMIT}. Here is
the result that establishes the non-collapsing in this case.



\begin{proposition}[Non-collapsing at scales r<=r_{i+1} follows directly from canonical neighborhoods]
\label{prop:noncollapsing-small-scale-easy-case}
\uses{def:strong-canonical-neighborhood-assumption}
Let $x\in {\mathcal M}$ with ${\bf t}(x)=t$ and with $R(x)=r^{-2}\ge r_{i+1}^{-2}$. Then there is $\kappa>0$ depending only
on $C$  such that ${\mathcal M}$ is $\kappa$-non-collapsed at $x$; that is to say,
if $R(x)=r^{-2}$ with $r\le r_{i+1}$, then $\Vol\,B(x,t,r)\ge \kappa r^3$, or $x$ is contained in a component of $M_t$ with positive sectional curvature.
\end{proposition}

\begin{proof}
 Since
$R(x)\ge r_{i+1}^{-2}$, by assumption any such $x$ has a strong
$(C,\epsilon)$-canonical neighborhood. If this neighborhood is a
strong $\epsilon$-neck centered at $x$ then the result is clear for
a non-collapsing constant $\kappa$ which is universal. If the
neighborhood is an $\epsilon$-round component containing $x$, then
$x$ is contained in a component of positive sectional curvature. Likewise,
if $x$ is contained in a $C$-component then by definition it is contained in a component of its
time-slice with positive sectional curvature.


Lastly, we suppose that $x$ is contained in the core $Y$ of a
$(C,\epsilon)$-cap ${\mathcal C}$. Let $r'>0$ be such that the
supremum of $|\Rm|$ on $B(x,t,r')$ is $(r')^{-2}$.
Then, by the
definition
 of a $(C,\epsilon)$-cap,
$\vol\,B(x,t,r')\ge C^{-1}(r')^3$.
Clearly, $r'\le r$ and there is a point $y\in \overline{B(x,t,r')}$
with $R(y)=(r')^{-2}$. On the other hand, by the
definition of a $(C,\epsilon)$-cap, we have $R(y)/R(x)\le C$, so that $r'/r\ge C^{-1/2}$.
Thus, the
volume of $B(x,t,r)$ is at least $C^{-5/2}r^3$.

This completes an examination of all cases and establishes the
proposition.
\end{proof}



\section{Minimizing ${\mathcal L}$-geodesics exist when
$R(x)\le r^{-2}_{i+1}$: the statement}


The proof of the non-collapsing result when $R(x)=r^{-2}$ with
$r_{i+1}<r\le \epsilon$ is much more delicate. As we indicated
above, it is analogous to the proof of non-collapsing for Ricci
flows given in Chapter~\ref{noncoll}. That is to say, in this case
the result is proved using the length function on the Ricci flow
with surgery and the monotonicity of the reduced volume. Of course,
unlike the case of Ricci flows treated in Chapter~\ref{noncoll},
here not all points of a Ricci flow with surgery ${\mathcal M}$ can
be reached by minimizing ${\mathcal L}$-geodesics, or rather more
precisely by minimizing ${\mathcal L}$-geodesics contained in the
open subset of smooth points of ${\mathcal M}$. (It is only for the
latter ${\mathcal L}$-geodesics that the analytic results of
Chapter~\ref{lengthfn} apply.) Thus, the main thing to establish in
order to prove non-collapsing is that for any Ricci flow with
surgery $({\mathcal M},G)$ satisfying the hypothesis of
Proposition~\ref{KAPPALIMIT} there are minimizing ${\mathcal
L}$-geodesics in the open subset of smooth points of ${\mathcal M}$
to `enough' of ${\mathcal M}$ so that we can run the same reduced
volume argument that worked in Chapter~\ref{noncoll}. Here is the
statement that tells us that there are minimizing ${\mathcal
L}$-geodesics to `enough' of ${\mathcal M}$.





\begin{proposition}[Existence and compactness of minimizing L-geodesics avoiding the surgeries]
\label{prop:minimizing-L-geodesics-exist}
\uses{prop:kappa-limit-noncollapsing}\label{delta0ri+1}
For each $r_{i+1}$ with $0<r_{i+1}\le r_i$, there is
$\delta=\delta(r_{i+1})>0$ (depending implicitly on $t_0$, $C$,
$\epsilon$, and $i$) such that if $\overline\delta(t)\le \delta$ for
all $t\in [T_{i-1},T_{i+1}]$ then the following holds. Let
$({\mathcal M},G)$ be a Ricci flow with surgery satisfying the
hypothesis of Proposition~\ref{KAPPALIMIT} with respect to the given
sequences and $r_{i+1}$, and let $x\in {\mathcal M}$ have ${\bf
t}(x)=T$ with $T\in [T_i,T_{i+1})$. Suppose that for some $r\ge
r_{i+1}$ the parabolic neighborhood $P(x,r,T,-r^2)$ exists in
${\mathcal M}$ and $|\Rm|\le r^{-2}$ on this neighborhood. Then
there is an open subset $U$ of ${\bf t}^{-1}[T_{i-1},T)$ contained
in the open subset of smooth manifold points  of ${\mathcal M}$ with
the following properties:
\begin{enumerate}
\item[(1)] For every $y$ in $U$ there is a minimizing ${\mathcal L}$-geodesic
connecting $x$ to $y$.
\item[(2)] $U_t=U\cap {\bf t}^{-1}(t)$ is non-empty for every $t\in [T_{i-1},T)$
\item[(3)] For each $t\in [T_{i-1},T)$ the restriction of ${\mathcal
L}$ to $U_t$ achieves its minimum and that minimum is at most
$3\sqrt{(T-t)}$.
\item[(4)] The subset of $U$ consisting of all $y$ with the property that
${\mathcal L}(y)\le {\mathcal L}(y')$ for all $y'\in {\bf
t}^{-1}({\bf t}(y))$ has the property that its intersection with
${\bf t}^{-1}(I)$ is compact for every compact interval $I\subset
[T_{i-1},T)$.
\end{enumerate}
\end{proposition}


The basic idea in proving this result is to show that all paths
beginning at $x$ and parameterized by backward time that come close
to the exposed regions have large ${\mathcal L}$-length. If we can
establish this, then the existence of such paths will not be an
impediment to using the analytic estimates from
Chapter~\ref{lengthfn} to show that for each $t\in [T_{i-1},T)$
there is a point whose ${\mathcal L}$-length from $x$ is at most
$3\sqrt{T-t}$, and that the set of points that minimize the
${\mathcal L}$-length from $x$ in a given time-slice form a compact
set.

Given Proposition~\ref{delta0ri+1},  arguments from
Chapter~\ref{noncoll} will be applied to complete the proof of
Proposition~\ref{KAPPALIMIT}.




\section{Evolution of neighborhoods of surgery caps}


 We begin this analysis required to prove Proposition~\ref{delta0ri+1}
 by studying the
evolution of surgery caps.
Proposition~\ref{altern} below is the main result along these lines.
Qualitatively, it says that if the surgery control parameter
$\delta$ is sufficiently small, then as a surgery cap evolves in a
Ricci flow with surgery it stays near the rescaled version of the
standard flow for any rescaled time less than one unless the entire
cap is removed (all at once) by some later surgery. In that case,
the evolution of the cap is close to the rescaled version of the
standard flow until it is removed. Using this result we will show
that if a path parameterized by backward time has final point  near
a surgery cap and has initial point with scalar curvature not too
large, then this path must enter this evolving neighborhood either
from the `top' or `size' and because of the estimates that we derive
in this chapter such a path  must have large ${\mathcal L}$-length.


\begin{proposition}[Evolution of a surgery cap stays close to the standard flow until removed]
\label{prop:surgery-cap-evolution-alternative}
\uses{def:surgery-cap,def:curvature-pinched-toward-positive-surgery}\label{altern}
Given $A<\infty$, $\delta''>0$ and $0<\theta<1$, there is
$\delta''_0=\delta''_0(A,\theta,\delta'')$ ($\delta_0''$ also
depends on $r_{i+1}$, $C$, and $\epsilon$, which are all now fixed)
such that the following holds. Suppose that $({\mathcal M},G)$ is a
Ricci flow with surgery defined for $0\le t<T$ with surgery control
parameter $\overline\delta(t)$. Suppose that it satisfies the strong
$(C,\epsilon)$-canonical neighborhood assumption at all points $x$
with $R(x)\ge r_{i+1}^{-2}$. Suppose also that $({\mathcal M},G)$ has
curvature that is pinched toward positive. Suppose that there is a
surgery at some time $\bar t$ with $T_{i-1}\le \bar t<T$ with $\bar
h$ as the surgery scale parameter. Set $T'=\mathrm{min}(T,\bar t+\theta
\bar h^2)$. Let $p\in M_{\bar t}$ be the tip of the cap of a surgery
disk. Then, provided that $\bar\delta(\bar t)\le \delta''_0$ one of
the following holds:
\begin{enumerate}
\item[(a)] There is an embedding $\rho\colon B(p,\bar t,A\bar h)\times [\bar t,T')\to {\mathcal M}$
compatible with time and the vector field. Let $g'(t),\ \bar t\le
t<T'$, be the one-parameter family of metrics on $B(p,\bar t,A\bar
h)$ given by $\rho^*G$. Shifting this family by $-\bar t$ to make
the initial time $0$ and scaling it by $(\bar h)^{-2}$ produces a
family of metrics $g(t),\ 0\le t<\mathrm{min}((T-\bar t)\bar
h^{-2},\theta)$, on $B_{g}(p,0,A)$ that are within $\delta''$ in the
$C^{[1/\delta'']}$-topology of the standard flow on the ball of
radius $A$ at time $0$ centered at the tip of its cap.
\item[(b)] There is $\bar t_+\in (\bar t,T')$ and an embedding
$B(p,\bar t,A\bar h)\times [\bar t,\bar t_+)\to {\mathcal M}$
compatible with time and the vector field so that the previous item
holds with $\bar t_+$ replacing $T'$. Furthermore, for any $t<\bar
t_+$ but sufficiently close to $\bar t_+$ the image of $B(p,\bar
t,A\bar h)\times\{t\}$ is contained in the region $D_{t}\subset M_t$
that disappears at time $\bar t_+$.
\end{enumerate}

\end{proposition}




\begin{proof}
\uses{prop:standard-flow-asymptotic-cylinder}
The method of proof is to assume that the result is false and take
a sequence of counterexamples with surgery control parameters
$\delta_n$ tending to zero. In order to derive a contradiction we
need to be able to take smooth limits of rescaled versions of these
Ricci flows with surgery, where the base points are the tips of the
surgery caps. This is somewhat delicate since the surgery cap is not
the result of moving forward for a fixed amount of time under Ricci
flow, and consequently Shi's theorem does not apply. Fortunately,
the metrics on the cap are bounded in the $C^\infty$-topology so
that Shi's theorem with derivatives does apply. Let us start
by examining limits of the sort we need to take.

\begin{claim}[Uniform derivative bounds for the rescaled surgery-cap metric]
\label{claim:surgery-cap-curvature-bound-uniform}
Let $(N,g_N)$ be a strong $\delta'$-neck with $N_0$ its middle half.
Suppose that $({\mathcal S},g)$ is the result of doing surgery on
(the central $2$-sphere) of $N$, adding a surgery cap ${\mathcal C}$
to $N^-$. Let $h$ be the scale of $N$. Let $({\mathcal S}_0(N),g')$
be the union of $N^-_0\cup {\mathcal C}$ with its induced metric as
given in Section~\ref{necksurgery}, and let $({\mathcal
S}_0(N),\widehat g_0)$ be the result of rescaling $g_0$ by $h^{-2}$.
Then for every $\ell<\infty$ there is a uniform bound to
$|\nabla^\ell \Rm_{\widehat g_0}(x)|$ for all $x\in {\mathcal
S}_0(N)$.
\end{claim}

\begin{proof}
\uses{thm:shi-derivative-estimates}
Since $(N,g_N)$ is a strong $\delta'$-neck of scale $h$, there is a
Ricci flow on $N$ defined for backward time $h^2$. After rescaling
by $h^{-2}$ we have a flow defined for backward time $1$.
Furthermore, the curvature of the rescaled flow is bounded on the
interval $(-1,0]$. Since the closure of $N_0$ in $N$ is compact, the
restriction of $h^{-2}g_N$ to $N_0\subset N$ at time $0$ is
uniformly bounded in the $C^\infty$-topology by Shi's theorem
(Theorem~\ref{shi}). The bound on the $k^{th}$-derivatives of the
curvature depends only on the curvature bound and hence can be taken
to be independent of $\delta'>0$ sufficiently small and also
independent of the strong $\delta'$-neck $N$. Gluing in the cap with
a  $C^\infty$-metric that converges smoothly to the standard initial
metric $g_0$ as $\delta'$ tends to zero using a fixed
$C^\infty$-partition of unity produces a family of manifolds
uniformly bounded in the $C^\infty$-topology.
\end{proof}

This leads immediately to:

\begin{corollary}[Smooth limit of rescaled surgery caps is the standard initial metric]
\label{cor:surgery-cap-smooth-limit}
Given a sequence of $\delta'_n\rightarrow 0$ and strong
$\delta'_n$-necks $(N(n),g_{N(n)})$ of scales $h_n$ and results of
surgery $(S_0(N(n)),g(n))$ with tips $p_n$ as in the previous claim,
then after passing to a subsequence there is a smooth limit
$(S_\infty,g_\infty,p_\infty)$ of a subsequence of the
$(S_0(N(n)),h_n^{-2}g_0(n)),p_n)$. This limit is the metric from Section~\ref{necksurgery}
that gives the standard initial conditions for a surgery cap.
\end{corollary}

\begin{proof}
That there is a smooth limit of a subsequence is immediate from the
previous claim. Since the $\delta_n$ tend to zero, it is clear that the
limiting metric is the standard initial metric.
\end{proof}

\begin{lemma}[Rescaled Ricci flows with surgery near a surgery cap converge to the standard flow]
\label{lem:rescaled-surgery-limit-standard-flow}
Suppose that we have a sequence of $3$-dimensional Ricci flows
with surgeries $({\mathcal M}_n,G_n)$ that satisfy the strong
$(C,\epsilon)$-canonical neighborhood assumption with parameter
$r_{i+1}$, and have curvature pinched toward positive. Suppose that
there are surgeries in ${\mathcal M}_n$ at times $t_n$ with surgery
control parameters $\delta'_n$ and scales $h_n$. Let $p_n$ be the
tip of a surgery cap for the surgery at time $ t_n$. Also suppose
that there is $0\le\theta_n<1$ such that for every $A<\infty$,
for all $n$ sufficiently large
there are embeddings
$B(p_a, t_n,A h_n)\times [t_n, t_n+ h^2_n\theta_n)\to {\mathcal
M}_n$ compatible with time and the vector field. Suppose that
$\delta'_n\rightarrow 0$ and $\theta_n\rightarrow \theta<1$ as $n\rightarrow
\infty$.  Let $({\mathcal M}'_n,G'_n,p_n)$ be the Ricci flow with
surgery obtained by shifting time by $-t_n$ so that surgery occurs
at ${\bf t}=0$ and rescaling by $h_n^{-2}$ so that the scale of the
surgery becomes one. Then, after passing to a subsequence, the
sequence converges smoothly to a limiting flow
$(M_\infty,g_\infty(t),(p_\infty,0)),\ 0\le t<\theta$. This limiting flow is isomorphic to
the restriction of the standard flow to time $0\le t<\theta$.
\end{lemma}


\begin{proof}
\uses{thm:shi-with-initial-derivatives}
Let $Q<\infty$ be an upper bound for the scalar curvature of the
standard flow on the time interval $[0,\theta)$. Since
$\delta'_n\rightarrow 0$, according to the previous corollary, there is
a smooth limit at time $0$ for a subsequence, and this limit is the
standard initial metric. Suppose that, for some $0\le
\theta'<\theta$, we have established that there is a smooth limiting
flow on $[0,\theta']$. Since the initial conditions are the standard
solution, it follows from the uniqueness statement in
Theorem~\ref{stdsoln} that in fact the limiting flow is
isomorphic to the restriction of the standard flow to this time
interval. Then the scalar curvature of the limiting flow is bounded
by $Q$. Hence, for any $A<\infty$, for all $n$ sufficiently large,
the scalar curvature of the restriction of $G'_n$ to the image of
$B_{G'_n}(p_n,0,2A)\times[0,\theta']$ is bounded by $2Q$. According
to Lemma~\ref{controlnbhd} there is an $\eta>0$ and a constant
$Q'<\infty$, each depending only on $Q$, $r_{i+1}$, $C$ and
$\epsilon$, such that for all $n$ sufficiently large, the scalar
curvature of the restriction of $G_n'$ to $B_{G_n'}(p_n,0,A)\times
[0,\mathrm{min}(\theta'+\eta,\theta_n))$ is bounded by $Q'$. Because of
the fact that the curvature is pinched toward positive, this implies
that on the same set the sectional curvatures are uniformly bounded.
Hence, by  Shi's theorem with derivatives
(Theorem~\ref{shiw/deriv}), it follows that there are uniform bounds
for the curvature in the $C^\infty$-topology. Thus, passing to a
subsequence we can extend the smooth limit to the time interval
$[0,\theta'+\eta/2]$ unless $\theta'+\eta/2\ge \theta$. Since $\eta$
depends on $\theta$ (through $Q$), but is independent of $\theta'$,
we can repeat this process extending the time-interval of definition
of the limiting flow by $\eta/2$ until $\theta'+\eta/2\ge \theta$.
Now suppose that $\theta'+\eta/2\ge \theta$. Then the argument shows
that by passing to a subsequence we can extend the limit to any
compact subinterval of $[0,\theta)$. Taking a diagonal sequence
allows us to extend it to all of $[0,\theta)$. By the uniqueness of
the standard flow, this limit is the standard flow.
\end{proof}

\begin{corollary}[Surgery-cap evolution is close to the standard flow up to time theta]
\label{cor:surgery-cap-close-to-standard-up-to-theta}\label{goodtotheta}
With the notation and assumptions of the previous lemma, for all
$A<\infty$, and any $\delta''>0$, then  for all $n$ sufficiently
large, the restriction of $G_n'$ to the image
$B_{G_n'}(p_n,0,A)\times [0,\theta_n)$ is within $\delta''$ in the
$C^{[1/\delta'']}$-topology of the restriction of the standard
solution to the ball of radius $A$ about the tip for time $0\le
t<\theta_n$.
\end{corollary}


\begin{proof}
\uses{thm:shi-with-initial-derivatives}
Let $\eta>0$ depending on $\theta$ (though $Q$) as well as
$r_{i+1}$, $C$ and $\epsilon$ be as in the proof of the previous
lemma, and take $0<\eta'<\eta$. For all $n$ sufficiently large
$\theta_n>\theta-\eta'$, and consequently for all $n$ sufficiently
large there is an embedding $B_{G_n}(p_n,t_n,Ah_n)\times
[t_n,t_n+h_n^2(\theta-\eta')]$ into ${\mathcal M}_n$ compatible with
time and with the vector field. For all $n$ sufficiently large, we
consider the restriction of $G_n'$ to $B_{G_n'}(p_n,0,A)\times
[0,\theta-\eta']$. These converge smoothly to the restriction of the
standard flow to the ball of radius $A$ on the time interval
$[0,\theta-\eta']$. In particular, for all $n$ sufficiently large,
the restrictions to these time intervals are within $\delta''$ in
the $C^{[1/\delta'']}$-topology of the standard flow. Also, for all $n$ sufficiently
large, $\theta_n-(\theta-\eta')<\eta$. Thus, by
Lemma~\ref{controlnbhd}, we see that the scalar curvature of $G_n'$
is uniformly bounded (independent of $n$) on
$B_{G_n'}(p_n,0,A)\times [0,\theta_n)$. By the assumption that the
curvature is pinched toward positive, this means that the sectional
curvatures of the $G_n'$ are also uniformly bounded on these sets,
and hence so are the Ricci curvatures. (Notice that these bounds are
independent of $\eta'>0$.) By Shi's theorem with derivatives
(Theorem~\ref{shiw/deriv}), we see that there are uniform bounds on
the curvatures in the $C^\infty$-topology on these subsets, and
hence bounds in the $C^\infty$-topology on the Ricci curvature.
These bounds are independent of both $n$ and $\eta'$. Thus, choosing
$\eta'$ sufficiently close to zero, so that $\theta_n-\eta'$ is also
close to $\theta$ for all $n$ sufficiently large, we see that for all such large $n$
and all
$t\in [\theta-\eta',\theta)$,  the
restriction of $G_n'$ to $B_{G_n'}(p_n,0,A)\times \{t\}$ is
arbitrarily close in the $C^{[1/\delta'']}$-topology to
$G_n'(\theta-\eta')$. The same is of course true of the standard
flow. This completes the proof of the corollary.
\end{proof}

Now we turn to the proof proper of Proposition~\ref{altern}. We fix
$A<\infty$, $\delta''>0$ and $\theta<1$. We are free to make $A$
larger so we can assume by Proposition~\ref{item5} that for the
standard flow the restriction of the flow to $B(p_0,0,A)\setminus
B(p_0,0,A/2)$ remains close to a standard evolving $S^2\times
[A/2,A]$ for time $[0,\theta]$. Let $K<\infty$ be a constant with
the property that $R(x,t)\le K$ for all $x\in B(p_0,0,A)$ in the
standard flow and all $t\in [0,\theta]$. If there is no
$\delta''_0>0$ as required, then we can find a sequence
$\delta'_n\rightarrow 0$ as $n\rightarrow\infty$ and Ricci flows with
surgery $({\mathcal M}_n,G_n)$ with surgeries at time $t_n$ with
surgery control parameter $\delta_n(t_n)\le \delta_n'$  and surgery
scale parameter $ h_n=h(r_{i+1}\delta_n( t_n),\delta_n(t_n))$
satisfying the hypothesis of the lemma but not the conclusion. Let
$T'_n$ be the final time of $({\mathcal M}_n,G_n)$. Let $\theta_n\le
\theta$ be maximal subject to the condition that there is an
embedding $\rho_n\colon B_{G_n}(x, t_n,Ah_n)\times [ t_n,
t_n+h_n^2\theta_n)\to {\mathcal M}_n$ compatible with time and the
vector field. Let $G'_n$ be the result of shifting the time by
$-t_n$ and scaling the result by $h_n^{-2}$. According to
Lemma~\ref{goodtotheta}, for all $n$ sufficiently large, the
restriction of $G_n'$ to the image of $\rho_n$ is within $\delta''$
in the $C^{[1/\delta'']}$-topology of the standard flow restricted
to the ball of radius $A$ about the tip of the standard solution on
the time interval $[0,\theta_n)$. If $\theta_n=\mathrm{min}(\theta,(T_n'-t_n)/h_n^2)$, then the first conclusion of
Proposition~\ref{altern} holds for $({\mathcal M}_n,G_n)$ for all
$n$ sufficiently large which contradicts our assumption that the
conclusion of this proposition holds for none of the $({\mathcal
M}_n,G_n)$. If on the other hand $\theta_n<\mathrm{min}(\theta,(T_n'-t_n)/h_n^2)$, we need only show that all of
$B(x_n,t_n,Ah_n)$ disappears at time $t_n+h_n^2\theta_n$ in order to
show that the second conclusion of Proposition~\ref{altern} holds
provided that $n$ is sufficiently large. Again this would contradict
the fact that the conclusion of this proposition holds for none of
the $({\mathcal M}_n,G_n)$.

So now let us suppose that $\theta_n<\mathrm{min}(\theta,(T'_n-t_n)/h_n^2)$. Since there is no further extension
in forward time for $B(p_n,t_n,A h_n)$, it must be the case that $
t_n+ h_n^2\theta_n$ is a surgery time and there is some flow line
starting at a  point of $B(p_n,t_n,A h_n)$ that does not continue to
time $t_n+h_n^2\theta_n$. It remains to show that in this case that
for any $t< t_n+ h_n^2\theta_n$ sufficiently close to $t_n+
h_n^2\theta_n$ we have $\rho_n\left(B_{G_n}(x, t_n,Ah_n)\times
\{t\}\right)\subset D_t$, the region in $M_t$ that disappears at
time $t_n+h_n^2\theta_n$.


\begin{claim}[Surgery spheres stay disjoint from the ball around the surgery cap]
\label{claim:ball-disjoint-from-surgery-spheres}
Suppose that $\theta_n<\mathrm{min}(\theta,(T'_n- t_n)/ h_n^2)$. Let
$\Sigma_{1},\ldots,\Sigma_{k}$ be the $2$-spheres along which
we do surgery at time $ t_n+h_n^2\theta_n$. Then for any $t< t_n+
h_n^2\theta_n$ sufficiently close to $ t_n+ h_n^2\theta_n$ the
following holds provided that $\delta_n'$ is sufficiently small. The image
$$\rho_n\left(B_{g_n}(x,t_n,A h_n)\times
\{t\}\right)$$ is disjoint from the images $\{\Sigma_{i}(t)\}$ of the $\{\Sigma_{i}\}$
under the
backward flow to time $t$ of the spheres $\Sigma_{i}$ along which
we do surgery at time $t_n+h_n^2\theta_n$.
\end{claim}

\begin{proof}
\uses{cor:surgery-cap-close-to-standard-up-to-theta,cor:embedded-sphere-isotopic-to-neck-fiber}
There is a constant $K'<\infty$ depending on $\theta$ such that for
the standard flow we have $R(x,t)\le K'$ for all $x\in B(p_0,0,A)$
and all $t\in [0,\theta)$ for the standard solution. Consider the
embedding $\rho_n\left(B(p_n,t_n,Ah_n)\times [t_n,t_n+h_n^2
\theta_n)\right)$. After time shifting by $-t_n$ and rescaling by
$h_n^{-2}$, the flow $G_n'$ on the image of $\rho_n$ is within
$\delta''$ of the standard flow. Thus, we see that for all $n$ sufficiently
large and for every point $x$ in the image of $\rho_n$ we have $
R_{G_n'}(x)\le 2K'$ and hence $R_{G_n}(x)\le 2K'h_n^{-2}$.


Let $h_n'$ be the scale of the surgery at time $t_n+ h_n^2\theta_n$.
(Recall that $h_n$ is the scale of the surgery at time $t_n$.)
Suppose that $\rho_n(B(p_n,t_n,Ah_n)\times \{t'\})$ meets one of the
surgery $2$-spheres $\Sigma_{i}(t')$ at time $t'$ at a point
$y(t')$. Then, for all $t\in [t',t_n+h_n^2\theta_n)$ we have
 the image $y(t)$ of $y(t')$ under the flow. All these points
$y(t)$ are points of intersection of $\rho_n(B(p,t_n,Ah_n)\times
\{t\})$ with $\Sigma_{i}(t)$. Since $y(t)\in \rho_n(B(p,t_n,A
h_n)\times\{t\})$, we have $ R(y(t))\le 2K' h_n^{-2}$. On the other
hand $R(y(t)) (h'_n)^2$ is within $O(\delta)$ of  $1$ as $t$ tends
to $t_n+h_n^2\theta_n$. This means that $h_n/h_n'\le \sqrt{3K'}$ for
all $n$ sufficiently large. Since the standard solution has
non-negative curvature, the metric is a decreasing function of $t$,
and hence the diameter of $B(p_0,t,A)$ is at most $ 2A$ in the standard
solution. Using Corollary~\ref{goodtotheta} we see that  for all $n$
sufficiently large, the diameter of $\rho_n\left(B(p,t_n,Ah_n)\times
\{t\}\right)$ is at most $Ah_n\le 4\sqrt{K'} Ah_n'$.
 This
means that for $\delta'_n$ sufficiently small the distance at time
$t$ from $\Sigma_{i}(t)$ to the complement of the $t$ time-slice
of the strong $\delta_n( t_n+ h_n^2\theta_n)$-neck $N_{i}(t)$
centered at $\Sigma_{i}(t)$ (which is at least
$(\delta'_n)^{-1}h_n'/2$) is much larger than the diameter of
$$\rho_n(B(p_n,t_n,Ah_n)\times \{t\}).$$  Consequently,  for all $n$
sufficiently large, the image $\rho_n(B(p_n,t_n,Ah_n)\times \{t\})$
is contained in  $N_{i}(t)$. But by our choice of $A$, and
Corollary~\ref{goodtotheta} there is an $\epsilon$-neck of rescaled
diameter approximately $Ah_n/2$ contained in $\rho_n(B(p_n,t_n,A
h_n)\times \{t\})$. By Corollary~\ref{s2isotopic} the spheres coming
from the neck structure in $$\rho_n(B(p_n,t_n,Ah_n)\times \{t\})$$ are
 isotopic in $N_{i}(t)$ to the central $2$-sphere of this neck.
This is a contradiction because in $N_{i}(t)$ the central
$2$-sphere is homotopically non-trivial whereas  the spheres in
$\rho_n(B(p_n,t_n,A\bar h_n)\times \{t\})$  clearly bound
$3$-disks.
\end{proof}

Since $\rho_n(B(p_n,t_n,Ah_n)\times\{t\})$ is disjoint from the
backward flow to time $t$ of all the surgery $2$-spheres
$\Sigma_{i}(t)$ and since $\rho_n(B(p_n,t_n,Ah_n)\times \{t\})$ is
connected,
 if there is a flow line starting at some point $z\in
B(p,t_n,Ah_n)$ that disappears at time $t_n+ h_n^2\theta_n$, then
the flow from every point of $B(p,t_n,Ah_n)$ disappears at time $
t_n+ h_n^2\theta_n$. This shows that if $\theta_n<\mathrm{min}(\theta,T'_n- t_n/h_n^{2})$, and if there is no extension of
$\rho_n$ to an embedding defined at time $t_n+h^2_n\theta_n$, then
all forward flow lines beginning at points of $B(p,t_n,A h_n)$
disappear at time $ t_n+h_n^2\theta_n$, which of course means that
for all $t<t_n+h^2_n\theta_n$ sufficiently close to
$t_n+h_n^2\theta_n$ the entire image $\rho_n(B(p,t_n,Ah_n)\times
\{t\})$ is contained in the disappearing region $D_t$. This shows
that for all $n$ sufficiently large, the second conclusion of
Proposition~\ref{altern} holds,  giving a contradiction.

This completes the proof of Proposition~\ref{altern}.
\end{proof}


\begin{remark}[A surgery-cap neighborhood may later be entirely removed by a further surgery]
\label{rem:ball-removed-by-later-surgery}
\uses{prop:surgery-cap-evolution-alternative}
Notice that it is indeed possible that $B_{G}(x,t,Ah)$ is removed at
some later time, for example as part of a capped  $\epsilon$-horn
associated to some later surgery.
\end{remark}

\section{A length estimate}

We use the result in the previous section about the evolution of
surgery caps to establish the length estimate on paths parameterized
by backward time approaching a surgery cap from above.

\begin{definition}[The constant overline delta_0 controlling closeness to the standard metric]
\label{def:overline-delta-0}\label{defnoverlinedelta0}
\uses{prop:standard-solution-scalar-curvature-lower-bound}
Let $c>0$ be the constant from Proposition~\ref{Restim}. Fix
$0<\overline\delta_0<1/4$ such that if $g$ is within
$\overline\delta_0$ of $g_0$ in the $C^{[1/\overline
\delta]}$-topology then $|R_{g'}(x)-R_{g_0}(x)|<c/2$ and $|\Ric_{g'}-\Ric_{g_0}|<1/4$.
\end{definition}


Here is the length estimate.


\begin{proposition}[Paths of bounded energy near a surgery cap have long integral length]
\label{prop:long-L-length-near-surgery-cap}
\uses{prop:surgery-cap-evolution-alternative}\label{bigell}
For any $\ell<\infty$ there is $A_0=A_0(\ell)<\infty$,
$0<\theta_0=\theta_0(\ell)<1$, and for any $A\ge A_0$ for the
constant
$\delta''=\delta''(A)=\delta''_0(A,\theta_0,\overline\delta_0)>0$
from Proposition~\ref{altern} the following holds. Suppose that
$({\mathcal M},G)$ is a Ricci flow with surgery defined for $0\le
t< T<\infty$. Suppose that it satisfies the strong
$(C,\epsilon)$-canonical neighborhood assumption at all points $x$
with $R(x)\ge r_{i+1}^{-2}$. Suppose also that the solution has
curvature pinched toward positive. Suppose that there is a surgery
at some time $\bar t$ with $T_{i-1}\le \bar t<T$ with
$\overline\delta(\bar t)$ as the surgery control parameter and with
$h$ as the surgery scale parameter. Then the following holds
provided that $\overline\delta(\bar t)\le \delta''$. Set $T'=\mathrm{min}(T,\bar t+h^2\theta_0 )$. Let $p\in M_{\bar t}$ be the tip of
the cap of a surgery disk at time $\bar t$. Suppose that $P(p,\bar
t,Ah,T'-\bar t)$ exists in ${\mathcal M}$. Suppose that we have
$t'\in [\bar t,\bar t+h^2/2]$ with $t'\le T'$, and suppose that we
have a curve $\gamma(\tau)$ parameterized by backward time $\tau\in
[0,T'- t']$ so that $\gamma(\tau)\in M_{T'-\tau}$ for all $\tau\in
[0,T'- t']$. Suppose that the image of $\gamma$ is contained in the
closure of $P(p,\bar t,Ah,T'-\bar t)\subset {\mathcal M}$. Suppose
further:
\begin{enumerate}
\item[(1)] either that $T'=\bar t +\theta_0 h^2\le T$ or
that $\gamma(0)\subset \partial B(p,\bar t,Ah)\times\{T'\}$; and
\item[(2)]  $\gamma(T'-t')\in B(p,\bar t,Ah/2)\times{t'}$.
\end{enumerate}
Then
$$\int_0^{T'- t'}\left(R(\gamma(t))+|X_\gamma(t)|^2\right)dt>\ell.$$

\end{proposition}




\begin{proof}
\uses{prop:surgery-cap-evolution-alternative,prop:standard-solution-scalar-curvature-lower-bound}
The logic of the proof is as follows. We
fix $\ell<\infty$. We shall determine the relevant value of
$\theta_0$ and then of $A_0$ in the course of the argument. Then for any $A\ge A_0$
we define
 $\delta''(A)= \delta''_0(A,\theta_0,\overline\delta_0)$, as in
Proposition~\ref{altern}.

The integral expression is invariant under time translation and also
under rescaling. Thus, we can (and do) assume that $\bar t=0$ and
that the scale $h$ of the surgery is $1$. We use the embedding of
$P(p,0,A,T')\to {\mathcal M}$ and write the restriction of the flow
to this subset as a one-parameter family of metrics $g(t),\ 0\le
t\le T'$, on $B(p,0,A)$. With this renormalization, $0\le t'\le
1/2$, also $T'\le \theta_0$,  and $\tau=T'-t$.


Let us first consider the case when  $T'=\theta_0\le T$. Consider the
standard flow $(\Ar^3,g_0(t))$, and let $p_0$ be its tip. According
to Proposition~\ref{Restim},  for all $x\in \Ar^3$ and all $t\in
[0,1)$ we have $R_{g_0}(x,t)\ge c/(1-t)$. By Lemma~\ref{altern} and
since we are assuming that $\overline\delta(\bar t)\le \delta''=
\delta''_0(A,\theta_0,\overline\delta_0)$, we have that $R(a,t)\ge
c/2(1-t)$ for all $a\in B(p,0,A)$ and all $t\in [0,\theta]$. Thus,
 we have
\begin{eqnarray*}
\int_0^{\theta_0-t'}\left(R(\gamma(\tau))+|X_\gamma(\tau)|^2\right)d\tau
& \ge & \int_{t'}^{\theta_0}
\frac{c}{2(1-t)}dt \\
& = & \frac{-c}{2}\left(\mathrm{log}(1-\theta_0)-\mathrm{log}(1-t')\right)dt \\
& \ge &  \frac{-c}{2}\left(\mathrm{log}(1-\theta_0)+\mathrm{log}(2)\right).
\end{eqnarray*}
 Hence, if
$\theta_0<1$ sufficiently close to $1$, the  integral will be $> \ell$.
This fixes the value of $\theta_0$.

\begin{claim}[Standard solution neck structure far from the tip persists up to time theta_0]
\label{claim:standard-solution-neck-at-infinity}
There is $A'_0<\infty$ with the property that for any $A\ge A'_0$
the restriction of the standard solution $g_0(t)$ to
$\left(B(p_0,0,A)\setminus B(p_0,0,A/2)\right)\times [0,\theta_0]$
is close to an evolving family $(S^2\times [A/2,A],h_0(t)\times
ds^2)$. In particular, for any $t\in [0,\theta_0]$, the
$g_0$-distance at time $t$ from $B(p_0,0,A/2)$ to the complement of
$B(p_0,0,A)$ in the standard solution is more than $A/4$.
\end{claim}

\begin{proof}
\uses{prop:standard-flow-asymptotic-cylinder}
This is immediate from Proposition~\ref{item5} and the fact that
$\theta_0<1$.
\end{proof}

Now fix $A_0=\mathrm{max}(A_0', 10\sqrt{\ell})$ and let $A\ge A_0$.


Since $\overline\delta_0<1/4$ and since $T'\le \theta_0$, for $\overline\delta(\bar t)\le
\delta''_0(A,\theta_0,\overline\delta_0)$
by Proposition~\ref{altern}
the $g(T')$-distance  between $B(p,0,A/2)$ and $\partial
B(p,0,A)$  is at least $A/5$.

 Since the
flow on $B(p,0,A)\times[0,T']$ is within $\overline\delta_0$ of
the standard solution, and since the curvature of the standard
solution is non-negative, for any horizontal tangent vector $X$ at
any point of $B(p,0,A)\times [0,T']$ we have that
$$\Ric_g(X,X)\ge -\frac{1}{4}|X|_{g_0}^2\ge -\frac{1}{2}|X|^2_g,$$ and hence
$$\frac{d}{dt}|X|_g^2\le |X|_g^2.$$
Because  $T'\le 1$, we see that
$$|X|^2_{g(T')}\le e\cdot|X|^2_{g(t)}<3|X|^2_{g(t)}$$
for any $t\in [0,T']$.

Now suppose that $\gamma(0)\in \partial B(p,0,A)\times\{T'\}$. Since
the image of $\gamma$ is contained in the closure of $P(p,0,A,T')$
for every $\tau\in [0,T']$ we have
$\sqrt{3}|X_\gamma(\tau)|_{g(T'-\tau)}\ge |X_\gamma(\tau)|_{g(T')}$.
Since the flow $g(t)$ on $P(p,0,A, T')$ is within
$\overline\delta_0$ in the $C^{[1/\overline\delta_0]}$-topology of
the standard flow on the corresponding parabolic neighborhood,
$R(\gamma(t))\ge 0$ for all $t\in [0,T']$. Thus, because of these
two estimates we have
\begin{equation}\label{*}
\int_0^{T'-t'}\left(R(\gamma(\tau))+|X_\gamma(\tau)|^2\right)d\tau\ge
\int_0^{T'-t'}
\frac{1}{3}|X_\gamma(\tau)|_{g(T')}^2d\tau.\end{equation}
 Since $\gamma(0)\in
\partial B(p,0,A)\times \{T'\}$ and $\gamma(T')\in B(p,0,A/2)$, it follows from
Cauchy-Schwarz that
\begin{eqnarray*}(T'-t')^2\int_0^{T'}|X_\gamma(\tau)|_{g(T')}^2d\tau &
\ge & \left(\int_0^{T'-t'}|X_\gamma(\tau)|_{g(T')}d\tau\right)^2
\\
& \ge & \left(d_{g(T')}(B(p,0,A/2),\partial
B(p,0,A))\right)^2\ge \frac{A^2}{25}. \end{eqnarray*}
 Since $T'-t'<1$, it
immediately follows from this and Equation~(\ref{*}) that
$$\int_0^{T'-t'}\left(R(\gamma(\tau))+|X_\gamma(\tau)|^2\right)d\tau\ge
\frac{A^2}{75}.$$ Since $A\ge A_0\ge 10\sqrt{ \ell}$, this
expression is $>\ell$. This completes the proof of
Proposition~\ref{bigell}
\end{proof}








\subsection{Paths with short ${\mathcal L}_+$-length avoid the surgery caps}

Here we show that a path parameterized by backward time that ends in
a surgery cap (or comes close to it)  must have
long ${\mathcal L}$-length. Let $({\mathcal M},G)$ be a Ricci flow
with surgery, and let $x\in {\mathcal M}$ be a point with ${\bf
t}(x)=T\in (T_i,T_{i+1}]$. We suppose that these data satisfy the
hypothesis of Proposition~\ref{delta0ri+1} with respect to the given
sequences and $r\ge r_{i+1}>0$. In particular, the parabolic
neighborhood $P(x,T,r,-r^2)$ exists in ${\mathcal M}$ and $|\Rm|$ is
bounded on this parabolic neighborhood by $r^{-2}$.

Actually, here we do not work directly with the length function
${\mathcal L}$ defined from $x$, but rather with a closely related
function. We set $R_+(y)=\mathrm{max}(R(y),0)$.


\begin{lemma}[Curves with small L_+-length avoid neighborhoods of surgery caps]
\label{lem:small-L-plus-avoids-surgery-caps}
\uses{prop:long-L-length-near-surgery-cap}\label{L_0}
Given $L_0<\infty$, there is
$\overline\delta_1=\overline\delta_1(L_0,r_{i+1})>0$, independent of
$({\mathcal M},G)$ and $x$, such that if $\overline\delta(t)\le
\overline\delta_1$ for all $t\in [T_{i-1},T)$, then for any curve
$\gamma(\tau),\ 0\le \tau\le \tau_0$, with $\tau_0\le T-T_{i-1}$,
parameterized by backward time with $\gamma(0)=x$ and with
$${\mathcal L}_+(\gamma)=\int_0^{\tau_0}\sqrt{\tau}
\left(R_+(\gamma(\tau))+|X_\gamma|^2\right)d\tau< L_0$$ the
following two statements hold: \begin{enumerate} \item[(2)] Set
$$\tau'=\mathrm{min}\left(\frac{r_{i+1}^4}{(256)L_0^2},\mathrm{ln}(\sqrt[3]{2})r_{i+1}^2\right).$$
Then for all $\tau\le \mathrm{min}(\tau',\tau_0)$ we have $\gamma(\tau)\in P(x,T,r/2,-r^2)$.
\item[(2)] Suppose that  $\bar t\in [T-\tau_0,T)$ is a surgery time with $p$
being the tip of the surgery cap at time $\bar t$ and with the scale
of the surgery being $\bar h$. Suppose   $t'\in [\bar t,\bar t+\bar
h^2/2]$ is such that there is an embedding
$$\rho\colon B(p,\bar t,(50+A_0)\bar h)\times [\bar t,t']\to {\mathcal M}$$
compatible with time and the vector field. Then the image of $\rho$
is disjoint from the image of $\gamma$. 
\end{enumerate}
\end{lemma}

\begin{remark}[Recollection: the surgery cap radius in terms of the surgery scale]
\label{rem:surgery-cap-radius-recollection}
\uses{def:surgery-cap}
Recall that $(A_0+4)\bar h$ is the radius of the surgery cap (measured in the
rescaled version of the standard initial metric) that is glued in when
performing surgery with scale $\bar h$.
\end{remark}






\begin{proof}
\uses{prop:surgery-cap-evolution-alternative,prop:long-L-length-near-surgery-cap}
We define $\ell=L_0/\sqrt{\tau'}$, then define $A=\mathrm{max}(A_0(\ell),2(50+A_0))$ and $\theta=\theta_0(\ell)$. Here,
$A_0(\ell)$ and $\theta_0(\ell)$ are the constants in
Proposition~\ref{bigell}.
  Lastly, we require $\overline\delta_1\le \delta''(A)$ from
Proposition~\ref{bigell}. Notice that, by construction,
 $\delta''(A)=\delta''_0(A,\theta,\overline\delta_0)$ from
Proposition~\ref{altern}. Thus, if $p$ is the tip of a surgery cap at time $\bar t$
with the scale of the surgery being $\bar h$, then it follows that
for any $\Delta t\le \theta$, if there is an embedding
$$\rho\colon B(p,\bar t,A\bar h)\times [\bar t,\bar t+\bar h^2\Delta t)\to {\mathcal M}$$
 compatible with time and the vector field, then the induced flow (after time shifting by $-\bar t$
 and scaling by $(\bar h)^{-2}$ is within $\overline\delta_0$ in the
 $C^{[1/\overline\delta_0]}$-topology of the standard solution.
 In particular, the scalar curvature
at any point of the image of $\rho$  is positive and is within a
multiplicative factor of two of the scalar curvature at the
corresponding point of the standard flow.

 Recall that we have $r\ge r_{i+1}$ and
that $P(x,T,r,-r^2)$ exists in ${\mathcal M}$ and that $|\Rm|\le
r^{-2}$ on this parabolic neighborhood.  We begin by proving by
contradiction that there is no $\tau\le  \tau'$ with the property
that $\gamma(\tau)\not\in P(x,T,r/2,-r^2)$. Suppose there is such a
$\tau\le \tau'$. Notice that by construction
$\tau'<r_{i+1}^2<r^2$. Hence,  for the first $\tau''$ with the
property that $\gamma(\tau'')\not\in P(x,T,r/2,-r^2)$ the point
$\gamma(\tau'')\in
\partial B(x,T,r/2)\times \{T-\tau''\}$.

\begin{claim}[Lower bound on the velocity integral before exiting a fixed ball]
\label{claim:energy-lower-bound-exit-ball}
$\int_0^{\tau''}|X_\gamma(\tau)|d\tau> r/2\sqrt{2}$.
\end{claim}

\begin{proof}
Since $|\Rm|\le r^{-2}$ on $P(x,T,r,-r^2)$, we have $|\Ric|\le
2r^{-2}$ on $P(x,T,r,-\tau'')$. Thus, for any tangent vector $v$ at
a point of $B(x,T,r)$ we have
$$\left|\frac{d(\langle v,v\rangle_{G(T-\tau)})}{d\tau}\right|\le 2r^{-2}\langle v,v\rangle_{G(T-\tau)}$$
for all
$\tau\in [0,\tau'']$. Integrating gives that for any $\tau\le
\tau''$ we have
$$\mathrm{exp}(-2r^{-2}\tau'')\langle v,v\rangle_{G(T)}\le \langle
v,v\rangle_{G(T-\tau)} \le \mathrm{exp}(2r^{-2}\tau'')\langle
v,v\rangle_{G(T)}.$$ Since  $\tau''\le \tau'$ and $r\ge r_{i+1}$ by
the assumption on $\tau'$ we have
$$\mathrm{exp}(2r^{-2}\tau'')\le \mathrm{exp}(2\sqrt[3]{2})< 2.$$ This
implies that for all $\tau\le \tau''$ we have
$$\frac{1}{\sqrt{2}}| X_\gamma(\tau)|_{G(T)}<
|X_\gamma(\tau)|_{G(T-\tau)}<\sqrt{2}|X_\gamma(\tau)|_{G(T)},$$
and hence
$$\int_0^{\tau''}|X_\gamma(\tau)|d\tau>
\frac{1}{\sqrt{2}}\int_0^{\tau''}|X_\gamma(\tau)|_{G(T)}\ge
\frac{r}{2\sqrt{2}},$$ where we use the fact that
$d_T(\gamma(0),\gamma(\tau''))=r/2$.
\end{proof}

Applying Cauchy-Schwarz to $\tau^{1/4}|X_\gamma|$ and $\tau^{-1/4}$ on the
interval $[0,\tau'']$ yields
\begin{eqnarray*}\int_0^{\tau''}\sqrt{\tau}\left(R_+(\gamma(\tau))+|X_\gamma(\tau)|^2\right)d\tau &
\ge & \int_0^{\tau''}\sqrt{\tau}|X_\gamma(\tau)|^2d\tau
\\
& \ge &
\frac{\left(\int_0^{\tau''}|X_\gamma(\tau)|d\tau\right)^2}{\int_0^{\tau''}\tau^{-1/2}d\tau}
\\ & > &
\frac{r^2}{16\sqrt{\tau''}}\ge L_0.
\end{eqnarray*}
Of course, the integral from $0$ to $\tau''$ is less than or equal
the entire integral from $0$ to $\tau_0$ since the integrand is
non-negative, contradicting the assumption that ${\mathcal
L}_+(\gamma)\le L_0$. This completes the proof of the first
numbered statement.


We turn now to the second statement. We impose a further condition
on $\overline\delta_1$. Namely, require that
$\overline\delta_1^2<r_{i+1}/2$. Since $r_{i}\le r_0\le \epsilon<1$,
we have $\overline\delta_1^2r_i<r_{i+1}/2$. Thus, the scale of the
surgery, $\bar h$, which is $\le \overline\delta_1^2r_i$ by
definition, will also be less than $r_{i+1}/2$, and hence there is
no point of $P(x,T,r,-r^2)$ (where the curvature is bounded by
$r^{-2}\le r_{i+1}^{-2}$) in the image of $\rho$ (where the scalar
curvature is greater than $(\bar h)^{-2}/2>2r_{i+1}^{-2}$). Thus, if
$\tau'\ge \tau_0$ we have completed the proof. Suppose that
$\tau'<\tau_0$. It suffices to establish that for every $\tau_1\in
[\tau',\tau_0]$ the point $\gamma(\tau_1)$ is not contained in the
image of $\rho$ for any surgery cap and any $t'$ as in the
statement. Suppose that in fact there is $\tau_1\in [\tau', \tau_0]$
with $\gamma(\tau_1)$ contained in the image of $\rho(B(p,\bar
t,(A_0+50)\bar h)\times [\bar t,t'])$ where $\bar t\le t'\le \bar
t+(\bar h)^2/2$ and where $p$ is the tip
 of some surgery cap at time $\bar t$. We estimate
\begin{eqnarray}\nonumber
\lefteqn{\int_0^{\tau_0}\sqrt{\tau}\left(R_+(\gamma(\tau))+|X_\gamma(\tau)|^2\right)d\tau} & & \\
& \ge &
\int_{\tau'}^{\tau_0}\sqrt{\tau}\left(R_+(\gamma(\tau))+|X_\gamma(\tau)|^2\right)d\tau
\nonumber
\\ & \ge &
\sqrt{\tau'}\int_{\tau'}^{\tau_1}\left(R_+(\gamma(\tau))+|X_\gamma(\tau)|^2\right)d\tau.
\label{tau_1est}
\end{eqnarray}

 Let $\Delta t\le
T-\bar t$ be the supremum of the set of $s$ for which there is a
parabolic neighborhood $P(p,\bar t,A\bar h,s)$ embedded in ${\bf
t}^{-1}((-\infty,T])\subset {\mathcal M}$. Let $\Delta t_1=\mathrm{min}(\theta \bar h^2,\Delta t)$. We consider $P(p,\bar t,A\bar
h,\Delta t_1)$. First, notice that since $\bar h\le
\overline\delta_1^2r_{i}<r_{i+1}/2$,
the scalar curvature on $P(p,\bar
t,A\bar h,\Delta t_1)$ is larger than $(\bar
h)^{-2}/2>r_{i+1}^{-2}\ge r^{-2}$. In particular, the parabolic
neighborhood $P(x,T,r,-r^2)$ is disjoint from $P(p,\bar t,A\bar
h,\Delta t_1)$. This means that there is some $\tau''\ge \tau'$ such
that $\gamma(\tau'')\in
\partial P(p,\bar t,A\bar h,\Delta t_1)$ and $\gamma|_{[\tau'',\tau_1]}\subset
P(p,\bar t,A\bar h,\Delta t_1)$. There are two cases to consider.
The first is when $\Delta t_1=\theta \bar h^2$, $\tau''=T-(\bar
t+\Delta t_1)$ and $\gamma(\tau'')\in B(p,\bar t,A\bar
h)\times\{\bar t+\Delta t_1\}$. Then, according to
Proposition~\ref{bigell},
\begin{equation}\label{**}
\int_{\tau''}^{\tau_1}R_+(\gamma(\tau))d\tau> \ell.\end{equation}

Now let us consider the other case. If $\Delta t_1<\theta \bar h^2$,
this means that either $\bar t+\Delta t_1=T$ or, according to
Proposition~\ref{altern}, at the time $\bar t+\Delta t_1$ there is a
surgery that removes all of $B(p,\bar t,A\bar h)$. Hence, under
either possibility it must be the case that $\gamma(\tau'')\in
\partial B(p,\bar t,A\bar h)\times \{T-\tau''\}$. Thus, the remaining
case to consider is when, whatever $\Delta t_1$ is,
$\gamma(\tau'')\subset
\partial B(p,\bar t,A\bar h)\times \{T-\tau''\}$.
Lemma~\ref{bigell} and the fact that $R\ge 0$ on $P(p,\bar t,A\bar
h,\Delta t_1)$ imply that
$$\ell<\int_{\tau''}^{\tau_1}\left(R(\gamma(\tau))+|X_\gamma(\tau)|^2\right)d\tau =
\int_{\tau''}^{\tau_1}\left(R_+(\gamma(\tau))+|X_\gamma(\tau)|^2\right)d\tau.$$

Since $\ell=L_0/\sqrt{\tau'}$ and $\tau''\ge \tau'$, it follows from
Equation~(\ref{tau_1est}) that in both
cases
$${\mathcal L}_+(\gamma)\ge
\int_{\tau''}^{\tau_1}\sqrt{\tau}\left(R_+(\gamma(\tau))+|X_\gamma(\tau)|^2\right)d\tau>
\ell\sqrt{\tau'}=L_0,$$ which contradicts our hypothesis. This
completes the proof of Lemma~\ref{L_0}.
\end{proof}

\subsection{Paths with small energy avoid
 the disappearing regions}


At this point we have shown that paths of small energy do not approach
the surgery caps from above. We also need to rule out that they can be arbitrarily
close from below. That is to say,
 we need to see that paths whose ${\mathcal L}$-length is not
too large avoid neighborhoods of the disappearing regions at all
times just before the surgery time at which they disappear. Unlike
the previous estimates which were universal for all $({\mathcal
M},G)$ satisfying the hypothesis of Proposition~\ref{delta0ri+1}, in
this case the estimates will depend on the Ricci flow with surgery.
First, let us fix some notation.

\begin{definition}[The disappearing region J near a surgery time]
\label{def:disappearing-region-J}
\uses{def:surgery-operation}\label{J}
 Suppose that $\bar t$ is a surgery
time, that $\tau_1>0$,  and that there are no other surgery times in
the interval $(\bar t-\tau_1,\bar t]$. Let $\{\Sigma_i(\bar t)\}_i$
be the $2$-spheres on which we do surgery at time $\bar t$. Each
$\Sigma_i$ is the central $2$-sphere of a strong $\delta$-neck
$N_i$. We can flow the cylinders $J_0(\bar
t)=\cup_is_{N_i}^{-1}(-25,0])$ backward to any time
 $t\in (\bar t-\tau_1,\bar t]$. Let $J_0(t)$ be
the result. There is an induced function, denoted
$\coprod_is_{N_i}(t)$, on $J_0(t)$. It takes values in $ (-25,0]$.
We denote the boundary of $J_0(t)$ by $\coprod_i\Sigma_i(t)$. Of
course, this boundary is the result of flowing
$\coprod_i\Sigma_i(\bar t)$ backward to time $t$. (These backward
flows are possible since there are no surgery times in $(\bar
t-\tau_1,\bar t)$.) For each $t\in [\bar t-\tau_1,\bar t)$ we also
have the disappearing region $D_t$: -- the region that disappears at
time $\bar t$. It is an open submanifold whose boundary is
$\coprod_i\Sigma_i(t)$. Thus, for every $t\in (\bar t-\tau_1,\bar
t)$ the subset $J(t)=J_0(t)\cup D_t$ is an open subset of $M_t$. We
define
$$J(\bar t-\tau_1,\bar t)=\cup_{t\in (\bar t-\tau_1,\bar t)}J(t).$$
Then $J(\bar t-\tau_1,\bar t)$ is an open subset of ${\mathcal M}$.

\end{definition}










\begin{lemma}[Paths of bounded kinetic energy avoid the disappearing region]
\label{lem:short-energy-avoids-disappearing-region}
\uses{def:disappearing-region-J}\label{15.17}
Fix a Ricci flow with surgery $({\mathcal M},G)$, a point $x\in
{\mathcal M}$ and constants $r\ge r_{i+1}>0$ as in the statement of
Proposition~\ref{delta0ri+1}. For any $1<\ell<\infty$ the following
holds. Suppose that $\bar t\in [T_{i-1},T)$ is a surgery time and
that  $\gamma(\tau)$ is a path with $\gamma(\tau)\in M_{\bar
t-\tau}$. Let $\{p_1,\ldots,p_k\}$ be the tips of all the surgery
caps at time $\bar t$ and let $\bar h$ be the scale of surgery at
time $\bar t$. Suppose that for some $0<\tau_1\le \ell^{-1} \bar
h^2$ there are no surgery times in the interval $(\bar t-\tau_1,\bar
t)$. We identify all $M_t$ for $t\in [\bar t-\tau_0,\bar t)$ with
$M_{\bar t-\tau_1}$ using the flow. Suppose that $\gamma(0)\in
M_{\bar t}\setminus \cup_{i=1}^kB(p_i,\bar t,(50+A_0)\bar h)$, and
lastly, suppose that
$$\int_{0}^{\tau_1}|X_\gamma(\tau)|^2d\tau\le \ell.$$
Then $\gamma$ is disjoint from the open subset $J(\bar t-\tau_1,\bar
t))$ of ${\mathcal M}$.
\end{lemma}




\begin{proof}
Suppose that the lemma is false and let $\gamma\colon [0,\bar
\tau]\to {\mathcal M}$ be a path satisfying the hypothesis of the
lemma with $\gamma(\bar \tau)\in J(\bar t-\tau_1,\bar t)$.
 Since $$\gamma(0)\in
M_{\bar t}\setminus \cup_iB(p_i,\bar t,(50+A_0)\bar h),$$ if follows
that $\gamma(0)$ is separated from the boundary of
$s_{N_i}^{-1}(-25,0]$ by distance at least $20\bar h$. Since the $
J_0(t)$ are contained in the disjoint union of strong $\delta$-necks
$N_i$ centered at the $2$-spheres along which we do surgery, and
since $\tau_1\le \bar h^2/\ell<\bar h^2$, it follows that, provided
that $\delta$ is sufficiently small, for every $t\in [\bar
t-\tau_1,\bar t)$, the metric on $J_0(t)$ is at least
$1/2$ the metric on $ J_0(\bar t)$.
 It follows that, for $\delta$ sufficiently small, if there is a $\tau\in [0,\tau_1]$
with $\gamma(\tau)\in J(t)$ then $\int_0^{\tau_1}|X_\gamma|d\tau> 10\bar h$.
Applying Cauchy-Schwarz we see that
$$\int_0^{\tau_1}|X_\gamma(\tau)|^2d\tau\ge (10\bar h)^2/\tau_1.$$
Since $\tau_1\le \ell^{-1}(\bar h)^2$, we see that
$$\int_0^{\tau'}|X_\gamma(\tau)|^2d\tau>\ell,$$
contradicting our hypothesis.
\end{proof}


\subsection{Limits of a sequence of paths with short ${\mathcal L}$-length}


Now using Lemmas~\ref{L_0} and~\ref{15.17} we show that it is possible to take limits
of certain types of sequences of paths parameterized by backward time to create
minimizing ${\mathcal L}$-geodesics.

We shall work with a compact subset of ${\bf t}^{-1}([T_{i-1},T])$ that is
obtained by removing appropriate open neighborhoods of the exposed regions.

\begin{definition}[Compact subset Y(ell) of smooth points avoiding all singular regions]
\label{def:compact-subset-Y-ell}
\uses{def:disappearing-region-J,prop:long-L-length-near-surgery-cap}
Fix $\ell<\infty$. Let $\theta_0=\theta_0(\ell)$ be as in
Proposition~\ref{bigell}. For each surgery time $\bar t\in
[T_{i-1},T]$, let $h(\bar t)$ be the scale of the surgery.
 Let $p_1,\ldots,p_k$ be
the tips of the surgery caps at time $\bar t$. For each $1\le j\le k$,
we consider $B_j(\bar
t)=B(p_j,\bar t,(A_0+10)h(\bar t))$, and we let $\Delta t_j\le \mathrm{min}(\theta_0,(T-\bar t)/h^2(\bar t))$ be maximal subject to the condition that
there is an embedding $\rho_j\colon B_j(\bar t)\times [\bar t,\bar t+h^2(\bar
t)\Delta t_j)$ into ${\mathcal M}$ compatible with time and the vector field.
Clearly, $B'_j=B(p_j,\bar t,(10+A_0)h)\cap C_{\bar t}$ is contained in $J(\bar
t)$. Let $\bar t'$ be the previous surgery time if there is one, otherwise set
$\bar t'=0$. Also for each $\bar t$ we set $\tau_1(\ell,\bar t)=\mathrm{min}\left(h(\bar t)^2/\ell,\bar t-\bar t'\right)$. For each $t\in (\bar
t-\tau_1(\ell,\bar t),\bar t)$ let $\widetilde J(t)\subset J(t)$ be the union
of $D_t$, the disappearing region at time $t$, and $\coprod_iB_i'(t)$,  the
result of flowing $\coprod_iB'_i$ backward to time $t$. Then we set
$\widetilde J(\bar t-\tau_1(\ell,\bar t),\bar t)\subset J(\bar
t-\tau_1(\ell,\bar t),\bar t)$ equal to the union over $t\in (\bar
t-\tau_1(\ell,\bar t),\bar t)$ of $\widetilde J(t)$.





 By construction,
for each surgery time $\bar t$, the union  $$\nu_\mathrm{sing}(\ell,\bar
t)=\widetilde J(\bar t-\tau_1(\ell,\bar t),\bar t)\cup\bigcup
_iB_i\times [\bar t,\bar t+h^2(\bar t)\Delta t_i)$$ is an open
subset of ${\mathcal M}$ containing all the exposed regions and
singular points at time $\bar t$.

We define $Y(\ell )\subset {\bf t}^{-1}([T_{i-1},T])$ to be the
complement of the $\cup_{\bar t}\nu_\mathrm{sing}(\ell,\bar t)$ where
the union is over all surgery times $\bar t\in [T_{i-1},T]$.
Clearly, $Y(\ell )$ is a closed subset of ${\bf
t}^{-1}([T_{i-1},T])$ and hence $Y(\ell )$ is a compact subset
contained in the open subset of smooth points of ${\mathcal M}$.
(Notice that $Y(\ell )$ depends on $\ell $ because $\tau_1(\ell,\bar
t)$ and $\theta_0$ depend on $\ell$.)
\end{definition}


\begin{proposition}[Existence, semicontinuity, and compactness of limits of short L-length paths]
\label{prop:L-geodesics-exist-and-limit}
\uses{lem:small-L-plus-avoids-surgery-caps}\label{Lgeoexist}
Fix $0<L<\infty$. Set
$$L_0=L+4(T_{i+1})^{3/2}.$$ Suppose that for all
$t\in [T_{i-1},T_{i+1}]$, the surgery control parameter
$\overline\delta(t)\le \overline\delta_1(L_0,r_{i+1})$
where the right-hand side is the constant from
Lemma~\ref{L_0}. Suppose that $\gamma_n$ is a sequence of paths in
$({\mathcal M},G)$ parameterized by backward time $\tau\in
[0,\bar\tau]$ with $\bar \tau\le T-T_{i-1}$, with $\gamma_n(0)=x$ and
with
$${\mathcal L}(\gamma_n)\le L$$
for all $n$. Then:
\begin{enumerate}
\item[(1)]
After passing to a subsequence, there is a limit $\gamma$ defined on
$[0,\bar\tau]$. The limit $\gamma$ is a continuous path and is a
uniform limit of the $\gamma_n$. The limit is contained in the open subset
of smooth points of ${\mathcal M}$ and has finite
${\mathcal L}$-length satisfying
$${\mathcal L}(\gamma)\le \mathrm{liminf}_{n\rightarrow\infty}{\mathcal L}(\gamma_n).$$
\item[(2)]
If there is a point $y\in M_{T-\bar\tau}$ such that
$\gamma_n(\bar\tau)=y$ for all $n$, and if the $\gamma_n$ are a
sequence of paths parameterized by backward time from $x$ to $y$
with $\mathrm{lim}_{n\rightarrow\infty} {\mathcal L}(\gamma_n)$ being no greater than
the ${\mathcal L}$-length of any path from $x$ to $y$, then the limit
$\gamma$ of a subsequence is a minimizing ${\mathcal L}$-geodesic
connecting $x$ to $y$ contained in the open subset of smooth points of ${\mathcal M}$.
\item[(3)] There is $\ell<\infty$ depending only on $L$ such that
any path $\gamma$ parameterized by backward time from $x$ to a
point $y\in {\bf t}^{-1}([T_{i-1},T))$ whose ${\mathcal L}$-length
is at most $L$ is contained in the compact subset $Y(\ell)$ given in
the previous definition.
\end{enumerate}
\end{proposition}


\begin{proof}
\uses{lem:small-L-plus-avoids-surgery-caps,prop:surgery-cap-evolution-alternative,prop:long-L-length-near-surgery-cap,lem:l-geodesic-existence}
Given $L_0$, we set
$$\tau'=\mathrm{min}\left(\frac{r_{i+1}^4}{(256)L_0^2},\mathrm{ln}(\sqrt[3]{2})r_{i+1}^2\right)$$
as in Lemma~\ref{L_0} and then define $\ell=L_0/\sqrt{\tau'}$. We
also let $A=\mathrm{min}(2(50+A_0),A_0(\ell))$ and
$\theta_0=\theta_0(\ell)$ as in Proposition~\ref{bigell}. Lastly, we
let
$\overline\delta_1(L_0,r_{i+1})=\delta''(A)=\delta''(A,\theta_0,\overline\delta_0)$
from Propositions~\ref{bigell} and~\ref{altern}. We suppose that
$\delta(t)\le \overline\delta_1(L_0,r_{i+1})$ for all $t\in
[T_{i-1},T]$.

Let $\bar t\in [T_{i-1},T]$ be a surgery time, and let $\bar h$ be
the scale of the surgery at this time. For each surgery cap
${\mathcal C}$ with tip $p$ at a time $\bar t\in [T_{i-1},T]$ let
$\Delta t({\mathcal C})$ be the supremum of those $s$ with $0\le s\le \theta_0\bar
h^2$ for which there is an embedding
$$\rho_{\mathcal C}\colon
B(p,\bar t,2(A_0+50)\bar h)\times[\bar t,\bar t+s)\to{\mathcal M}$$
compatible with time and the vector field. We set
$$P_0({\mathcal C})=\rho_{\mathcal C}\left(B(p,\bar t,(A+50)\bar h)\times
[\bar t,\bar t+\mathrm{min}(\bar h^2/2,\Delta t({\mathcal C})))\right).$$

\begin{claim}[A short path from x misses the extended surgery-cap neighborhood P_0]
\label{claim:path-misses-P0-cap-neighborhood} Any path $\gamma$ beginning at $x$ and parameterized by backward
time misses $P_0({\mathcal C})$ if ${\mathcal L}(\gamma)<L$.
\end{claim}

\begin{proof}
\uses{lem:small-L-plus-avoids-surgery-caps}
Set $\tau_0=T-\bar t$. Of course, $\tau_0\le T-T_{i-1}\le T_{i+1}-T_{i-1}$.
Consider the restriction of $\gamma$ to $[0,\tau_0]$.
We have
\begin{eqnarray*}
\lefteqn{\int_0^{\tau_0}\sqrt{\tau}\left(R_+(\gamma_n(\tau))+|X_{\gamma_n}(\tau)|^2\right)d\tau
}  \\
& \le &
\int_0^{T-T_{i-1}}\sqrt{\tau}\left(R_+(\gamma_n(\tau))+|X_{\gamma_n}(\tau)|^2\right)d\tau \\
& \le & \int_0^{T-T_{i-1}}\sqrt{\tau}\left(R(\gamma_n(\tau))+|X_{\gamma_n}(\tau)|^2\right)d\tau+
\int_0^{T-T_{i-1}}
6\sqrt{\tau}d\tau \\
& = &
\int_0^{T-T_{i-1}}\sqrt{\tau}\left(R(\gamma_n(\tau))+|X_{\gamma_n}(\tau)|^2\right)d\tau
+4(T-T_{i-1})^{3/2}
\\
& \le &
\int_0^{\tau_0}\sqrt{\tau}\left(R(\gamma_n(\tau))+|X_{\gamma_n}(\tau)|^2\right)d\tau+4(T_{i+1})^{3/2}
\end{eqnarray*}
Thus, the hypothesis that ${\mathcal L}(\gamma_n)\le L$ implies that
\begin{equation}\label{L'ineq}
\int_0^{\tau_0}\sqrt{\tau}\left(R_+(\gamma_n(\tau))+|X_{\gamma_n}(\tau)|^2\right)d\tau\le
L_0.\end{equation}
The claim now follows immediately from Lemma~\ref{L_0}.
\end{proof}

Now set $t'$ equal to the last surgery time before $\bar t$ or set $t'=0$ if $\bar t$ is
the first surgery time. We set $\tau_1(\bar t)$ equal to the minimum of $\bar t-t'$ and
$\bar h^2/\ell$.

Assume that $\gamma(0)=x$ and that ${\mathcal L}(\gamma)\le L$. It follows from
Lemma~\ref{L_0} that the restriction of the path $\gamma$ to $[0,\tau']$ lies
in a region where the Riemann curvature is bounded above by $r^{-2}\le
r_{i+1}^{-2}$. Hence, since $\bar h<\bar\delta(t)^2r_{i+1}\ll r_{i+1}$, this
part of the path  is disjoint from all strong $\delta$-necks (evolving
backward for rescaled time $(-1,0]$). That is to say, $\gamma|_{[0,\tau']}$ is
disjoint from $J_0(t)$ for every $t\in (\bar t-\tau_1(\bar t),\bar t)$ for any
surgery time $\bar t\le T$. It follows immediately that $\gamma|_{[0,\tau']}$
is disjoint from $J(\bar t-\tau_1(\bar t),\bar t)$.

\begin{claim}[A short path from x is disjoint from J near each surgery time]
\label{claim:path-disjoint-from-J-near-surgery}
For every surgery time $\bar t\in [T_{i-1},T]$, the path $\gamma$
starting at $x$ with ${\mathcal L}(\gamma)\le L$ is disjoint from
$J(\bar t,\bar t-\tau_1(\bar t))$.
\end{claim}

\begin{proof}
\uses{lem:short-energy-avoids-disappearing-region}
By the remarks above, it suffices to consider surgery times $\bar
t\le T-\tau'$. It follows immediately from the previous claim that
for any surgery time $\bar t$, with the scale of the surgery being
$\bar h$ and with $p$ being the tip of a surgery cap at this time,
we have $\gamma$ is disjoint from $B(p,\bar t,(50+A_0)\bar h)$.
Also,
$$\int_{T-\bar t}^{T-\bar t+\tau_1(\bar t)}\sqrt{\tau}|X_\gamma(\tau)|^2d\tau
\le {\mathcal L}_+(\gamma)\le L_0.$$
Since we can assume $T-\bar t\ge \tau'$ this implies that
$$\int_{T-\bar t}^{T-\bar t-\tau_1(\bar t)}|X_\gamma(\tau)|^2d\tau\le L_0/\sqrt{\tau'}=\ell.$$
The claim is now immediate from Lemma~\ref{15.17}.
\end{proof}

From these two claims we see immediately that $\gamma$ is contained in the
compact subset $Y(\ell)$ which is contained in the open subset of smooth points
of ${\mathcal M}$. This proves the third item in the statement of the
proposition. Now let us turn to the limit statements.

Take a sequence of paths $\gamma_n$ as in the statement of
Proposition~\ref{Lgeoexist}.
 By Lemma~\ref{L_0} the
restriction of each $\gamma_n$ to the interval $[0,\mathrm{min}(\bar\tau,\tau')]$ is contained in $P(x,T,r/2,-r^2)$. The
arguments in the proof of Lemma~\ref{geoexist} (which involve
changing variables to $s=\sqrt{\tau}$) show that, after passing to a
subsequence, the restrictions of the $\gamma_n$ to $[0,\mathrm{min}(\bar\tau,\tau')]$ converge uniformly to a path $\gamma$ defined
on the same interval. Furthermore,
\begin{equation*}
\int_0^{\mathrm{min}(\bar\tau,\tau')}\sqrt{\tau}|X_\gamma(\tau)|^2d\tau\le \mathrm{liminf}_{n\rightarrow\infty}\int_0^{\mathrm{min}(\bar\tau,\tau')}\sqrt{\tau}|X_{\gamma_n}(\tau)|^2d\tau,
\end{equation*}
so that
\begin{eqnarray}\label{gamliminf1}\lefteqn{\int_0^{\mathrm{min}(\bar\tau,\tau')}\sqrt{\tau}\left(R(\gamma(\tau)+|X_\gamma(\tau)|^2\right)d\tau\le} \\
& & \mathrm{liminf}_{n\rightarrow\infty}\int_0^{\mathrm{min}(\bar\tau,\tau')}\sqrt{\tau}\left(R(\gamma_n(\tau))+|X_{\gamma_n}(\tau)|^2\right)d\tau.\nonumber
\end{eqnarray}

If $\bar\tau\le \tau'$, then we have established the existence of a
limit as required. Suppose now that $\bar\tau>\tau'$. We turn our
attention to the paths $\gamma_n|_{[\tau',\bar\tau]}$. Let
$T_{i-1}<\bar t\le T-\tau'$ be either a surgery time or $T-\tau'$, and let $t'$ be the
maximum of the last surgery time before $\bar t$ and $T_{i-1}$. We consider the
restriction of the $\gamma_n$ to the interval $[T-\bar t,T-t']$. As
we have seen, these restrictions  are disjoint from $J(\bar
t-\tau_1(\bar t),\bar t)$ and also from the exposed region at time
$\bar t$, which is denoted $E(\bar t)$,  and from $J_0(\bar t)$. Let
$$Y={\bf t}^{-1}([T-t',T-\bar t])\setminus
\left(J(\bar t-\tau_1(\bar t),\bar t)\cup (E(\bar t)\cup J_0(\bar
t))\right).$$ This is a compact subset with the property that any
point $y\in Y$ is connected by a backward flow line lying entirely
in $Y$ to a point $y(t')$ contained in $M_{t'}$.


 Since $Y$ is compact there is a finite upper bound on the
Ricci curvature on $Y$, and hence to ${\mathcal L}_\chi(G)$ at any
point of $Y$. Since all backward flow lines from points of $Y$
extend all the way to $M_{t'}$, it follows that there is a constant
$C'$ such that
$$|X_{\gamma_n}(\tau)|_{G(t')}\le C'|X_{\gamma_n}(\tau)|_{G(t)}$$
for all $t\in [t',\bar t]$. Our hypothesis that the ${\mathcal
L}(\gamma_n)$ are uniformly bounded, the fact that the curvature is
pinched toward positive and the fact that there is a uniform bound
on the lengths of the $\tau$-intervals implies that the
$$\int_{T-\bar t}^{T-t'} \sqrt{\tau}|X_{\gamma_n}(\tau)|^2d\tau$$
are uniformly bounded. Because $T-\bar t$ is at least $\tau'>0$, it
follows that the $\int_{T-\bar t}^{T-t'}|X_{\gamma_n}|^2d\tau$ have
a uniform upper bound. This then implies that there is a constant
$C_1$ such that for all $n$ we have
$$\int_{T-\bar t}^{T-t'}|X_{\gamma_n}(\tau)|_{G(t')}^2d\tau\le
C_1.$$ Thus, after passing to a subsequence, the $\gamma_n$ converge
uniformly to a continuous $\gamma$ defined on $[T-\bar t,T-t']$.
Furthermore, we can arrange that the convergence is a weak
convergence in $W^{1,2}$. This means that $\gamma$ has a derivative
in $L^2$ and
$$\int_{T-\bar t}^{T-t'}|X_{\gamma}(\tau)|^2d\tau\le \mathrm{liminf}_{n\rightarrow
\infty}\int_{T-\bar t}^{T-t'}|X_{\gamma_n}(\tau)|^2d\tau.$$


Now we do this simultaneously for all $\bar t=T-\tau'$ and for all the finite number of surgery times
in $[T_{i-1},T-\tau']$. This
gives a limiting path $\gamma\colon [\tau',\bar\tau]\to {\mathcal M}$.
Putting together the above inequalities we see that the limit
satisfies
\begin{equation}\label{gamliminf2}
\int_{\tau'}^{\bar\tau}\sqrt{\tau}\left(R(\gamma(\tau))+|X_\gamma(\tau)|^2\right)d\tau\le
\mathrm{liminf}_{n\rightarrow\infty}\int_{\tau'}^{\bar\tau}\sqrt{\tau}\left(R(\gamma_n(\tau))+
|X_{\gamma_n}(\tau)|^2\right)d\tau.
\end{equation}
 Since we have already
arranged that there is a limit on $[0,\tau']$, this produces a
limiting path $\gamma\colon [0,\tau_0]\to {\mathcal M}$. By
Inequalities~\ref{gamliminf1} and~\ref{gamliminf2} we see that
$${\mathcal L}(\gamma)\le \mathrm{liminf}_{i\rightarrow \infty}{\mathcal
L}(\gamma_n).$$
The limit lies in the compact subset $Y(\ell)$ and hence is contained
in the open subset of smooth points of ${\mathcal M}$.
 This completes the proof of the first statement of the proposition.

Now suppose, in addition to the above, that all the $\gamma_n$
have the same endpoint $y\in M_{T-\tau_0}$ and that $\mathrm{lim}_{n\rightarrow\infty}{\mathcal L}(\gamma_n)$ is less than or equal
to the ${\mathcal L}$-length of any path parameterized by backward
time connecting $x$ to $y$. Let $\gamma$ be the limit of a
subsequence as constructed in the proof of the first part of this
result. Clearly, by what we have just established, $\gamma$ is a path
parameterized by backward time from $x$ to $y$ and ${\mathcal
L}(\gamma)\le \mathrm{lim}_{n\rightarrow\infty}{\mathcal L}(\gamma_n)$.
This means that $\gamma$ is a minimizing ${\mathcal L}$-geodesic
connecting $x$ to $y$, an ${\mathcal L}$-geodesic contained in the open subset of smooth
points of ${\mathcal M}$.

This completes the proof of the proposition.
\end{proof}


\begin{corollary}[Existence of a minimizing L-geodesic to an earlier time-slice]
\label{cor:minimizing-geodesic-to-earlier-slice}
\uses{lem:small-L-plus-avoids-surgery-caps,prop:L-geodesics-exist-and-limit}\label{bardelta1}
Given $L<\infty$, let
$\overline{\delta}_1=\overline\delta_1(L+4(T_{i+1}^{3/2},r_{i+1})$ be as given in Lemma~\ref{L_0}.
If
$\overline\delta(t)\le \overline\delta_1$ for all $t\in [T_{i-1},T_{i+1}]$,
then for any $x\in {\bf t}^{-1}([T_i,T_{i+1}))$ and for any $y\in
M_{T_{i-1}}$, if there is a  path $\gamma$ parameterized by
backward time connecting $x$ to $y$ with ${\mathcal L}(\gamma)\le L$,
then there is a minimizing ${\mathcal L}$-geodesic contained in the
open subset of smooth points of ${\mathcal M}$ connecting $x$ to
$y$.
\end{corollary}

\begin{proof}
Choose an ${\mathcal L}$-minimizing sequence of paths from $x$  to
$y$ and apply the previous proposition.
\end{proof}




\subsection{Completion of the proof of Proposition~{\ref{delta0ri+1}}}

Having found a compact subset of the open subset of smooth points of ${\mathcal M}$
that contains all paths parameterized by backward time whose ${\mathcal L}$-length is not
too large, we are in a position to prove Proposition~\ref{delta0ri+1}, which states
the existence of a minimizing ${\mathcal L}$-geodesics in ${\mathcal M}$ from $x$
and gives estimates on their ${\mathcal L}$-lengths.



\begin{proof}(of Proposition~\ref{delta0ri+1}).
\uses{cor:minimizing-geodesic-to-earlier-slice,prop:L-geodesics-exist-and-limit}
Fix $r\ge r_{i+1}>0$. Let $({\mathcal M},G)$ and $x\in {\mathcal M}$ be
as in the statement of Proposition~\ref{delta0ri+1}. We set
$L=8\sqrt{T_{i+1}}(1+T_{i+1})$, and we set
$$\delta=\mathrm{min}\left(\delta_i,\overline \delta_1(L+4(T_{i+1})^{3/2},r_{i+1})\right),$$
where $\overline\delta_1$ is  as given in Lemma~\ref{L_0}.
Suppose that $\overline\delta(t)\le
\delta$ for all $t\in [T_{i-1},T_{i+1})$. We set $U$ equal to the
subset of ${\bf t}^{-1}([T_{i-1},T))$ consisting of all points $y$
for which there is a path $\gamma$ from $x$ to $y$, parameterized by
backward time, with ${\mathcal L}(\gamma)<L$. For each $t\in
[T_{i-1},T)$ we set $U_t=U\cap M_t$. According to
Corollary~\ref{bardelta1} for any $y\in U$ there is a minimizing
${\mathcal L}$-geodesic connecting $x$ to $y$ and this geodesic lies
in the open subset of ${\mathcal M}$ consisting of all the smooth
points of ${\mathcal M}$; in particular, $y$ is a smooth
point of ${\mathcal M}$. Let ${\mathcal L}_x\colon U\to \Ar$ be the
function that assigns to each $y\in U$ the ${\mathcal L}$-length of
a minimizing ${\mathcal L}$-geodesic from $x$ to $y$. Of course,
${\mathcal L}_x(y)<L$ for all  $y\in U$. Now let us show that the
restriction of ${\mathcal L}_x$ to any time-slice $U_t\subset U$
achieves its minimum along a compact set. For this, let $y_n\in U_t$
be a minimizing sequence for ${\mathcal L}_x$ and for each $n$ let
$\gamma_n$ be a minimizing ${\mathcal L}$-geodesic connecting $x$ to
$y_n$. Since ${\mathcal L}(\gamma_n)<L$ for all $n$, according to
Proposition~\ref{Lgeoexist}, we can pass to a subsequence that
converges to a limit, $\gamma$, connecting $x$ to some point $y\in
M_t$ and ${\mathcal L}(\gamma)\le \mathrm{inf}_{n}{\mathcal L}(\gamma_n)<L$. Hence, $y\in
U_t$, and clearly ${\mathcal L}_x|_{U_t}$ achieves its minimum at
$y$. Exactly the same argument with $y_n$ being a sequence of points
at which ${\mathcal L}_x|_{U_t}$ achieves its minimum shows  that
the subset of $U_t$ at which ${\mathcal L}_x$ achieves its minimum
is a compact set.


We set $Z\subset U$ equal to the set of $y\in U$ such that
${\mathcal L}_x(y)\le {\mathcal L}_x(y')$ for all $y'\in  U_{{\bf
t}(y)}$.

\begin{claim}[Compactness of the minimal-L-length locus Z' in each time-slice]
\label{claim:Z-prime-compact}
The subset  $Z'=\{z\in Z|{\mathcal L}_x(z)\le L/2\}$ has the property
that for any compact interval $I\subset [T_{i-1},T)$ the
intersection ${\bf t}^{-1}(I)\cap Z'$ is compact.
\end{claim}

\begin{proof}
Fix a compact interval $I\subset [T_{i-1},T)$. Let $\{z_n\}$ be a
sequence in $Z'\cap {\bf t}^{-1}(I)$. By passing to a subsequence we
can assume that the sequence ${\bf t}(z_n)=t_n$ converges to some
$t\in I$, and that ${\mathcal L}_x(z_n)$ converges to some $D\le
L/2$. Since the surgery times are discrete, there is a neighborhood
$J$ of $t$ in $I$ such that the only possible surgery time in $J$ is
$t$ itself. By passing to a further subsequence if necessary, we can
assume that $t_n\in J$ for all $n$. Fix $n$. First, let us consider
the case when $t_n\ge t$. Let $\gamma_n$ be a minimizing ${\mathcal
L}$-geodesic from $x$ to $z_n$. Then we form the path $\widehat
\gamma_n$ which is the union of $\gamma_n$ followed by the flow line
for the vector field $\chi$ from the endpoint of $\gamma_n$ to
$M_t$. (This flowline exists since there is no surgery time in the
open interval $(t,t_n]$.) If $t_n<t$, then we set $\widehat
\gamma_n$ equal to the restriction of $\gamma_n$ to the interval
$[0,T-t]$. In either case let $\hat y_n\in M_t$ be the endpoint of
$\widehat\gamma_n$. Since $M_t$ is compact, by passing to a
subsequence we can arrange that the $\hat y_n$ converge to a point $y\in
M_t$. Clearly,
 $\mathrm{lim}_{n\rightarrow\infty }z_n=y$.

It is also the case that $\mathrm{lim}_{n\rightarrow\infty}{\mathcal L}(\widehat\gamma_n)=
\mathrm{lim}_{n\rightarrow\infty}{\mathcal L}(\gamma_n)=D\le L/2$. This
means that $y\in U$ and that ${\mathcal L}_x(y)\le D\le L/2$. Hence,
the greatest lower bound of the values of ${\mathcal L}_x$ on $U_t$
is at most $D\le L/2$, and consequently $Z'\cap U_t\not=\emptyset$.
Suppose that the minimum value of ${\mathcal L}_x$ on $U_t$ is
$D'<D$. Let $z\in U_t$ be a point where this minimum value is
realized, and let $\gamma$ be a minimizing ${\mathcal L}$-geodesic
from $x$ to $z$. Then by restricting $\gamma$ to subintervals
$[0,t-\mu]$ shows that the minimum value of ${\mathcal L}_x$ on
$U_{t+\mu}\le (D'+D)/2$ for all $\mu>0$ sufficiently small. Also,
extending $\gamma$ by adding a backward vertical flow line from $z$
shows that the minimum value of ${\mathcal L}_x$ on $U_{t-\mu}$ is
at most $(D'+D)/2$ for all $\mu>0$ sufficiently small. (Such a
vertical flow line backward in time exists since $z\in U$ and hence
$z$ is contained in the smooth part of ${\mathcal M}$.) This
contradicts the fact that limit of the minimum values of ${\mathcal
L}_x$ on $U_{t_n}$ converge to $D$ as $t_n$ converges to $t$. This
contradiction proves  that the minimum value of ${\mathcal L}_x$ on
$U_t$ is $D$, and consequently the point $y\in Z'$. This proves that
$Z'\cap {\bf t}^{-1}(I)$ is compact, establishing the claim.
\end{proof}

At this point we  have established that Properties (1),(2), and (4);
So it remains only to prove Property (3) of Proposition~\ref{delta0ri+1}.
 To do this we define the reduced length function
$l_x\colon U\to \Ar$ by
$$l_x(q)=\frac{{\mathcal L}_x(q)}{2\sqrt{T-{\bf t}(q)}} \ \ \mathrm{and} \ \
l_x^\mathrm{min}(\tau)=\mathrm{min}_{q\in M_t}l_x(q).$$ We consider the
subset ${\mathcal S}$ of $\tau'\in (0,T-T_{i-1}]$ with $l_x^\mathrm{min}(\tau)\le L/2$ for all $\tau\le \tau'$. Recall that by the
choice of $L$, we have $3\sqrt{T-T_{i-1}}< L/2$. Clearly, the
minimum value of $l_x$ on $U_{T-\tau}$ converges to $0$
as $\tau\rightarrow 0$, implying that this set is non-empty. Also, from
its definition, ${\mathcal S}$ is an interval with $0$ being one
endpoint.

\begin{lemma}[The minimum reduced length along the surviving interval is at most 3/2]
\label{lem:reduced-length-min-bound-3-2}\label{lvalue}
Let $l_x^\mathrm{min}(\tau')$ be the minimum value of $l_x$ on
$U_{T-\tau'}$. For any $\tau\in {\mathcal S}$ we have $l_x^\mathrm{min}(\tau)\le 3/2$.
\end{lemma}

\begin{proof}
\uses{cor:minimizing-geodesic-to-earlier-slice,prop:minimizing-L-geodesics-exist,cor:generalized-flow-n-half-bound}
Given that we have already established Properties 1,2 and 4 of
Proposition~\ref{delta0ri+1}, this is immediate from
Corollary~\ref{compactcoro}.
\end{proof}


Now let us establish that ${\mathcal S}=(0,T-T_{i-1}]$. As we remarked above,
${\mathcal S}$ is a non-empty interval with $0$ as one endpoint. Suppose that
it is of the form $(0,\tau]$ for some $\tau<T-T_{i-1}$. Then by the previous
claim, we have $l_x^\mathrm{min}(\tau)\le 3/2$ so that there is an ${\mathcal
L}$-geodesic $\gamma$ from $x$ to a point $y\in M_{T-\tau}$ with ${\mathcal
L}(\gamma)\le 3\sqrt{\tau}< L/2$.
 This implies that for all $\tau'>\tau$ but sufficiently close to
$\tau$, there is a point $y(\tau')\in U_{T-\tau'}$ with ${\mathcal
L}_x(y(\tau'))<L/2$. This shows that all $\tau'$ greater than and sufficiently close
to $\tau$ are contained in ${\mathcal S}$.
 This is a
contradiction of the assumption that ${\mathcal S}=(0,\tau]$.

Suppose now that ${\mathcal S}$ is of the form $(0,\tau)$, and set $t=T-\tau$.
Let $t_n\rightarrow t$ and $z_n\in Z'\cap U_{t'}$. The same argument as above shows
that for every $n$ we have ${\mathcal L}_x(z_n)\le 3\sqrt{T-t_n}$. For all $n$
sufficiently large, there are no surgery times in the interval $(t,t_n)$.
Hence, by passing to a subsequence, we can arrange that the $z_n$ converge to a
point $z\in M_t$. Clearly, $${\mathcal L}_x(z)\le \mathrm{limsup}_{n\rightarrow\infty}{\mathcal L}_x(z_n)\le 3\sqrt{T-t},$$ so that $\tau\in
{\mathcal S}$. This again contradicts the assumption that ${\mathcal
S}=(0,\tau)$.

The only other possibility is that the set of $\tau$ is $(0,T-T_{i-1}]$ and the
minimum value of ${\mathcal L}$ on $U_t$ is at most $3\sqrt{T-t}$ for all $t\in
[T_{i-1},T)$. This is exactly the third property stated in
Proposition~\ref{delta0ri+1}. This completes the proof of that proposition.
\end{proof}




\section{Completion of the proof of
Proposition~{\ref{KAPPALIMIT}}}

Now we are ready to establish Proposition~\ref{KAPPALIMIT}, the
non-collapsing result. We shall do this by finding a parabolic
neighborhood whose size, $r'$, depends only on $r_i$, $C$ and
$\epsilon$, on which the sectional curvature is bounded by
$(r')^{-2}$ and so that the ${\mathcal L}$-distance from $x$ to any
point of the final time-slice of this parabolic neighborhood is
bounded. Recall that in Section~\ref{rsmall} we established it when
$R(x)=r^{-2}$ with $r\le r_{i+1}<\epsilon$. Here we assume that
$r_{i+1}<r\le \epsilon$.  Fix $\delta=\delta(r_{i+1})$ from
Proposition~\ref{delta0ri+1} and set $L=8\sqrt{T_{i+1}}(1+T_{i+1})$.

First of all, in Claim~\ref{newkappa0r0t0} we have  seen that there
is $\kappa_0$ such that ${\bf t}^{-1}[0,T_1]$ is $\kappa_0$
non-collapsed on scales $\le \epsilon$. Thus, we may assume that
$i\ge 1$.

Recall that ${\bf t}(x)=T\in (T_{i},T_{i+1}]$. Let $\gamma$ be an
${\mathcal L}$-geodesic contained in the smooth part of ${\mathcal
M}$ from $x$ to a point in $ M_{T_{i-1}}$ with ${\mathcal
L}(\gamma)\le 3\sqrt{T-T_{i-1}}$. That such a $\gamma$ exists was
proved in Proposition~\ref{delta0ri+1}. We shall find a point $y$ on
this curve with  $R(y)\le 2r_i^{-2}$. Then  we find a backward
parabolic neighborhood centered at $y$ on which ${\mathcal L}$ is
bounded and so that the slices have volume bounded from below. Then
we can apply the results from Chapter~\ref{noncoll} to establish the
$\kappa$ non-collapsing.

\begin{claim}[Existence of a point of curvature below r_i^{-2} along the L-geodesic]
\label{claim:existence-of-tau0-low-curvature-point}\label{tau0cl}
There is $\tau_0$  with $\mathrm{max}(\epsilon^2,T-T_i)\le \tau_0\le
T-T_{i-1}-\epsilon^2$ such that $R(\gamma(\tau_0))< r_i^{-2}$.
\end{claim}


\begin{proof}
Let $T'=\mathrm{max}(\epsilon^{2},T-T_i)$ and let
$T''=T-T_{i-1}-\epsilon^2$, and suppose that $R(\gamma(\tau))\ge
r_i^{-2}$ for all $\tau\in [T',T'']$. Then we see that
$$\int_{T'}^{T''}\sqrt{\tau}\left(R(\gamma(\tau)+|X_\gamma(\tau)|^2\right)d\tau
\ge \frac{2}{3}r_i^{-2}\left((T'')^{3/2}-(T')^{3/2}\right).$$ Since
$R\ge -6$ because the curvature is pinched toward positive, we see that
\begin{eqnarray*}
{\mathcal L}(\gamma) & \ge & \frac{2}{3}r_i^{-2}\left((T'')^{3/2}-(T')^{3/2}\right)-\int_0^{T'}
6\sqrt{\tau}d\tau-\int_{T''}^{T-T_{i-1}}6\sqrt{\tau}d\tau \\
& = & \frac{2}{3}r_i^{-2}\left((T'')^{3/2}-(T')^{3/2}\right)-
4(T')^{3/2}-4\left((T-T_{i-1})^{3/2}-(T'')^{3/2}\right).
\end{eqnarray*}

\begin{claim}[Auxiliary time-interval estimates for the low-curvature point]
\label{claim:tau0-estimates} We have the following estimates:
\begin{eqnarray*}
(T'')^{3/2}-(T')^{3/2} & \ge &
\frac{1}{4}(T-T_{i-1})^{3/2} \\
4(T')^{3/2} & \le &  4(T-T_{i-1})^{3/2} \\
4\left((T-T_{i-1})^{3/2}-(T'')^{3/2}\right) & \le &
\frac{2t_0}{25}\sqrt{(T-T_{i-1})}.\end{eqnarray*}
\end{claim}

\begin{proof}
Since $T_i-T_{i-1}\ge t_0$ and $T\ge T_i$, we see that
$T''/(T-T_{i-1})\ge 0.9$. If $T'=T-T_i$, then since
$T<T_{i+1}=2T_i=4T_{i-1}$ we have $T'/(T-T_{i-1})\le 2/3$. If
$T'=\epsilon^2$, since $\epsilon^2\le t_0/50$, and $T-T_{i-1}\ge
t_0$, we see that $T'\le (T-T_{i-1})/50$. Thus, in both cases we
have $T'\le 2(T-T_{i-1})/3$. Since $(0.9)^{3/2}>0.85$ and
$(2/3)^{3/2}\le 0.6$, the first inequality follows.

The second inequality is clear since $T'<(T-T_{i-1})$.

The last inequality is clear from the fact that
$T''=T-T_{i-1}-\epsilon^2$ and $\epsilon\le \sqrt{t_0/50}$.
\end{proof}




 Putting these together yields
 $${\mathcal L}(\gamma)\ge
 \left[\left(\frac{1}{6}r_i^{-2}-4\right)\left(T-T_{i-1}\right)-\frac{2t_0}{25}\right]
 \sqrt{T-T_{i-1}}.$$

Since
$$r_i^{-2}\ge r_0^{-2}\ge \epsilon^{-2}\ge 50/t_0,$$
and $T-T_{i-1}\ge t_0$ we see that \begin{eqnarray*}{\mathcal
L}(\gamma) & \ge &
\left[\left(\frac{50}{6t_0}-4\right)t_0-\frac{2t_0}{25}\right]\sqrt{T-T_{i-1}}
\\
& \ge & (8-5t_0)\sqrt{T-T_{i-1}} \\
& \ge & 4\sqrt{T-T_{i-1}}.
\end{eqnarray*}

(The last inequality uses the fact that $t_0=2^{-5}$.) But this
contradicts the fact that ${\mathcal L}(\gamma)\le
3\sqrt{T-T_{i-1}}$.
\end{proof}




Now fix $\tau_0$ satisfying Claim~\ref{tau0cl}.
 Let $\gamma_1$ be the
restriction of $\gamma$ to the subinterval $[0,\tau_0]$, and let
$y=\gamma_1(\tau_0)$. Again using the fact that $R(\gamma(\tau))\ge
-6$ for all $\tau$, we see that
\begin{equation}\label{Lgamma1}
{\mathcal L}(\gamma_1)\le {\mathcal L}(\gamma)+4(T-T_{i-1})^{3/2}\le
3(T_{i+1})^{1/2}+4(T_{i+1})^{3/2}.\end{equation}


Set $t'=T-\tau_0$. Notice that from the definition we
have $t'\le T_i$. Consider $B=B(y,t',\frac{r_i}{2C})$, and
define $\Delta=\mathrm{min}(r_i^2/16C,\epsilon^2)$. According to
Lemma~\ref{controlnbhd} every point $z$ on a backward flow line
starting in $B$ and defined for time at most $ \Delta$  has the
property that $R(z)\le 2r_i^{-2}$. For any surgery time $\bar t$ in
$[t'-\Delta,t')\subset [T_{i-1},T)$ the scale $\bar h$ of the
surgery at time $\bar t$ is $\le \delta(\bar t)^2 r_i$, and hence
every point of the surgery cap has scalar curvature at least
$D^{-1}\delta(\bar t)^{-4}r_i^{-2}$. Since $\overline\delta(\bar t)\le
\overline \delta\le \delta_0\le \mathrm{min}(D^{-1},1/10)$, it follows
that every point of the surgery cap has curvature at least
$\delta_0^{-3}r_i^{-2}\ge 1000r_i^{-2}$. Thus,
 no point
$z$ as above can lie in a surgery cap. This means that the entire
backward parabolic neighborhood $P(y,t',\frac{r_i}{2C},-\Delta)$
exists in ${\mathcal M}$, and the scalar curvature is bounded by $2
r_i^{-2}$ on this backward parabolic neighborhood. Because of the
curvature pinching toward positive assumption, there is $C'<\infty$
depending only on $r_i$ and such that the Riemann curvature is
bounded by $C'$ on $P(y,t',\frac{r_i}{2C},-\Delta)$.


Consider the one-parameter family of metrics $g(\tau),\ 0\le \tau\le
\Delta$, on $B(y,t',\frac{r_i}{2C})$ obtained by restricting the
horizontal metric $G$ to the backward parabolic neighborhood.
There is $0<\Delta_1\le \Delta /2$
depending only on $C'$ such that for every $\tau\in [0,\Delta_1]$ and
every non-zero tangent vector $v$ at a point of
$B(y,t',\frac{r_i}{2C})$ we have
$$\frac{1}{2}\le \frac{|v|^2_{g(\tau)}}{|v|^2_{g(0)}}\le 2.$$
Set $\hat r=\mathrm{min}(\frac{r_i}{32C},\Delta_1/2)$, so that $\hat r$ depends
only on $r_i$, $C$, and $\epsilon$. Set $t''=t'-\Delta_1$. Clearly,
$B(y, t'',\hat r)\subset B(y,t',\frac{r_i}{2C})$ so that $B(y,t'',\hat r)\subset
P(y,t',\frac{r_i}{2C},-\Delta)$. Of course, it then follows that the parabolic neighborhood
$P(y,t'',\hat r,-\Delta_1)$ exists in ${\mathcal M}$ and
$$P(y,t'',\hat r,-\Delta_1)\subset P(y,t',\frac{r_i}{2 C},-\Delta),$$ so that the
Riemann curvature is bounded above by $C'$ on the parabolic
neighborhood $P(y,t'',\hat r,-\Delta_1)$. We set $r'=\mathrm{min}(\hat r,(C')^{-1/2},\sqrt{\Delta_1}/2)$, so that $r'$ depends only on
$r_i$, $C$, and $\epsilon$. Then  the parabolic neighborhood
$P(y,t'',r',-(r')^2)$ exists in ${\mathcal M}$ and $|\Rm|\le
(r')^{-2}$ on $P(y,t'',r',-(r')^2)$. Hence, by the inductive
non-collapsing assumption either $y$ is contained in a component of $M_{t''}$ of positive sectional curvature or
$$\Vol\,B(y,t'',r')\ge
\kappa_i(r')^3.$$

If $y$ is contained in a component of $M_{t''}$ of positive sectional curvature, then by
  Hamilton's result, Theorem~\ref{flowtoround},
under Ricci flow the component of $M_{t''}$  containing $y$ flows forward as
a family of  components of positive sectional curvature until it disappears. Since there is path moving backwards in time
from $x$ to $y$, this means that the original point $x$ is contained
in a component of its time-slice with positive sectional curvature.

Let us consider the other possibility when $\Vol\,B(y,t'',r')\ge
\kappa_i(r')^3$.
 For each $z\in B(y,t'',r')$ let
 $$\mu_z\colon [T-t',T-t'']\to  B(y,t',\frac{r_i}{2C})$$ be
 the $G(t')$-geodesic connecting $y$ to $z$. Of course
 $$|X_{\mu_z}(\tau)|_{G(t')}\le \frac{r'}{\Delta_1}$$
 for every $\tau\in [0,\Delta_1]$. Thus,
 $$|X_{\mu_z}(\tau)|_{G(T-\tau)}\le
 \frac{\sqrt{2}r'}{\Delta_1}$$ for all $\tau\in [T-t',T-t'']$.
 Now we let $\tilde \mu_z$ be the resulting path parameterized by backward time
 on the time-interval $[T-t',T-t'']$.
 We estimate
 \begin{eqnarray*}{\mathcal
 L}(\tilde \mu_z) & = &
 \int_{T-t'}^{T-t''}\sqrt{\tau}\left(R(\tilde\mu_z(\tau))+|X_{\tilde\mu_z}(\tau)|^2\right)d\tau
 \\
& \le & \sqrt{T-t''}\int_{T-t'}^{T-t''}\left(2 r_i^{-2}+\frac{2(r')^2}{\Delta_1^2}\right)d\tau
\\
& \le & \sqrt{T-t''}(2r_i^{-2}\Delta_1+\frac{1}{2}) \le
\sqrt{T}\left(\frac{1}{16C}+\frac{1}{2}\right).
\end{eqnarray*}
In passing to the last inequality we use the fact, from the
definitions that $r'\le \sqrt{\Delta_1}/2$ and $\Delta\le
r_i^2/16C$, whereas $\Delta_1\le \Delta/2$.

Since $C>1$, we see that
$${\mathcal L}(\tilde \mu_z)\le \sqrt{T}.$$
Putting this together with the estimate, Equation~(\ref{Lgamma1}),
for ${\mathcal L}(\gamma_1)$ tells us that for each $z\in
B(y,t'',r')$ we have
$${\mathcal L}(\gamma_1*\tilde\mu_z)\le
4(T_{i+1})^{1/2}+4(T_{i+1})^{3/2}\le L/2.$$ Hence, by
Proposition~\ref{delta0ri+1} and the choice of $L$, there is a
minimizing ${\mathcal L}$-geodesic from $x$ to each point of
$B(y,t'',r')$ of length $\le L/2$, and these geodesics are contained
in the smooth part of ${\mathcal M}$. In fact, by
Proposition~\ref{Lgeoexist} there is a compact subset $Y$ of
the open subset of smooth points of ${\mathcal M}$ that contains all
the minimizing ${\mathcal L}$-geodesics from $x$ to points of
$B(y,t'',r')$.

Then, by Corollary~\ref{fullmeasure} (see also,
Proposition~\ref{lips}), the intersection, $B'$, of ${\mathcal U}_x$
with $B(y,t'',r')$
 is an open subset of full measure in $B(y,t'',r')$.
Of course, $\Vol\,B'=\Vol\,B(y,t'',r')\ge \kappa_i(r')^3$
and the function $l_x$ is bounded by $L/2$ on $B'$. It now follows
from Theorem~\ref{THM} that there is $\kappa>0$ depending only on
$\kappa_i$, $r'$, $\epsilon$ and $L$ such that $x$ is $\kappa$
non-collapsed on scales $\le \epsilon$. Recall that $L$ depends only
on $T_{i+1}$, and $r'$ depends only on $r_i,C,C'$ and $\epsilon$,
whereas $C'$ depends only on $r_i$. Thus, in the final analysis,
$\kappa$ depends only on $\kappa_i$ and $r_i$ (and $C$ and
$\epsilon$ which are fixed). This entire analysis assumed that for
all $t\in [T_{i-1},T_{i+1})$ we have the inequality
$\overline\delta(t)\le
\overline\delta_1(L+4(t_{i+1})^{3/2},r_{i+1})$ as in
Lemma~\ref{L_0}. Since $L$ depends only on $i$ and $t_0$, this shows
that the upper bound for $\delta$ depends only on $r_{i+1}$ (and on
$i$, $t_0$, $C$, and $\epsilon$). This completes the proof of
Proposition~\ref{KAPPALIMIT}.


